% The CM value formula as the pipeline computes it: one statement for every coset and every index,
% with a proof that treats composite squarefree level and non-fundamental discriminant on the same
% footing as the prime, fundamental case.  Part I is new; Part II collects the local computations
% from paper/kappa0-proof-standalone.tex, which this note is meant to supersede.
%
% STATUS: a draft for Eran and Sachi to read.  Nothing here is agreed, and the code and tests keep
% describing the m = 0 terms as "argued, not proved" until it is.
\documentclass[11pt,reqno]{amsart}
\usepackage[margin=1.15in]{geometry}
\usepackage{amsmath,amssymb,amsthm,mathtools}
\usepackage{booktabs}
\usepackage[colorlinks=true,linkcolor=blue,citecolor=blue,urlcolor=blue]{hyperref}
\theoremstyle{plain}
\newtheorem{theorem}{Theorem}[section]
\newtheorem{proposition}[theorem]{Proposition}
\newtheorem{lemma}[theorem]{Lemma}
\newtheorem{corollary}[theorem]{Corollary}
\theoremstyle{definition}
\newtheorem{definition}[theorem]{Definition}
\newtheorem{remark}[theorem]{Remark}
\DeclareMathOperator{\Diff}{Diff}
\DeclareMathOperator{\ord}{ord}
\DeclareMathOperator{\disc}{disc}
\DeclareMathOperator{\tr}{tr}
\DeclareMathOperator{\nrd}{nrd}
\DeclareMathOperator{\supp}{supp}
\DeclareMathOperator{\CM}{CM}
\DeclareMathOperator{\Res}{Res}
\newcommand{\Q}{\mathbb{Q}}
\newcommand{\Z}{\mathbb{Z}}
\newcommand{\C}{\mathbb{C}}
\newcommand{\R}{\mathbb{R}}
\newcommand{\A}{\mathbb{A}}
\newcommand{\Ldual}{L^{\vee}}
\newcommand{\Lm}{L_{-}}
\newcommand{\Lp}{L_{+}}
\newcommand{\e}{\mathbf{e}}
\newcommand{\parentref}[1]{\textup{[#1]}}
\begin{document}
\title[The CM value formula on $X_0^D(N)$]{The CM value formula for Borcherds forms on $X_0^D(N)$: \\ the coefficients $\kappa$ for every coset and every index}
\author{}
\date{\today}
\maketitle

\noindent\textbf{Status.} This note states, in one place, the formula the pipeline uses for the
values of Borcherds forms at CM points of $X_0^D(N)/W_{D,N}$ --- the function
\texttt{SchoferFormula} and the routines it calls --- and proves it under the hypotheses the
pipeline actually meets: squarefree level $N$ prime to $D$, any negative discriminant $d$ at which
there is a CM point, the Guo--Yang input forms. The proof is organised so that the level, the
conductor of $d$, and the ramification of the primes of $DN$ in $\Q(\sqrt d)$ enter only through
the local lattices, and the statement about those lattices is a single formula
(Theorem~\ref{thm:kappa}) rather than a list of cases; the case list (Part~II) is what one needs to
\emph{evaluate} it, and is where the earlier notes
\parentref{\texttt{level-prime-kappa.tex}, \texttt{kappa0-proof-standalone.tex}} are absorbed.
It is a draft: the arguments have been reviewed by agents, not yet agreed by the authors.

\tableofcontents

\part{The statement and its proof}

\section{Setting}\label{sec:setting}

\subsection*{The lattice}
Let $B$ be the indefinite quaternion algebra over $\Q$ of discriminant $D$ and $\mathcal O\subset B$
an Eichler order of level $N$, with $N$ squarefree and prime to $D$. Put
\[
  L=\{\alpha\in\mathcal O:\ \tr\alpha=0\},\qquad \langle\alpha,\beta\rangle=\tr(\alpha\beta'),
  \qquad Q(\alpha)=\tfrac12\langle\alpha,\alpha\rangle=\nrd(\alpha),
\]
$\beta'$ the canonical involution. Then $L$ is an even lattice of signature $(1,2)$
\cite[\S3]{Yang}, $V=L\otimes\Q$, $\Ldual$ its dual, $|\Ldual/L|=2(DN)^2$, and the level of
$L$ (the exponent of the discriminant form) is
\[
  M=4\cdot(\text{odd part of }DN)=\begin{cases}4DN&DN\text{ odd},\\2DN&DN\text{ even}.\end{cases}
\]
Let $\rho_L$ be the Weil representation of the metaplectic group on $\C[\Ldual/L]$ in the
normalisation of \cite[\S3]{Yang} (Borcherds'), $e_\eta$ the standard basis, and $\e(x)=e^{2\pi ix}$.

\subsection*{The input forms}
Let $\chi_\theta$ be the theta multiplier and $f$ a weakly holomorphic scalar form of weight
$\tfrac12$ on $\Gamma_0(M)$ with $f(\gamma\tau)=\chi_\theta(\gamma)(c\tau+d)^{1/2}f(\tau)$, with
poles at most at the cusps $\infty$ and $0$ (``$\{\infty,0\}$-weakly holomorphic''), with integral
principal parts in the sense of \cite[Lemma~25]{GY}. The pipeline's inputs are the eta quotients of
\cite[Lemma~13]{Yang} satisfying these conditions. Following \cite[(7)]{Yang}, \cite[(6)]{GY},
\[
  F=F_f=\sum_{\gamma\in\Gamma_0(M)\backslash SL_2(\Z)}(f|_{1/2}\gamma)\,\rho_L(\gamma^{-1})e_0
   =\sum_{\eta\in\Ldual/L}\ \sum_{m\in\Q}c_\eta(m)q^m\,e_\eta
\]
is a weakly holomorphic vector-valued form of weight $\tfrac12$ and type $\rho_L$ with
$O^+_{L,F}=O^+_L$ \cite[Lemmas~10, 12]{Yang}, and its principal part is \cite[Lemma~24]{GY}: with
$P(f)=\sum_{m>0}c_\infty(-m)q^{-m}$ the principal part of $f$ at $\infty$ and
$\frac{M\e(1/8)}{\sqrt{|\Ldual/L|}}P(f|_{1/2}S)=\sum_{n>0}b_nq^{-n/M}$ at the cusp $0$,
\begin{equation}\label{eq:GY24}
  c_\eta(-m)=\begin{cases}c_\infty(-m)+b_{mM}[\,mM\in\Z\,]&\eta=0,\\
  b_{mM}\,[\,mM\in\Z,\ Q(\eta)\equiv m\bmod\Z\,]&\eta\neq0,\end{cases}\qquad m>0 .
\end{equation}
In particular $c_\eta(-m)$ for $\eta\neq0$ depends on $\eta$ only through $Q(\eta)\bmod\Z$, and
$c_{-\eta}(m)=c_\eta(m)$ for all $\eta,m$ \cite[Thm.~13.3]{Borcherds}. Let $\Psi_F$ be the
Borcherds form of $F$ \cite[Thm.~A]{Yang} and $\psi_F(\tau)=\Psi_F(z(\tau))$ the corresponding
meromorphic modular form of weight $c_0(0)$ on $X=X_0^D(N)/W_{D,N}$ \cite[Lemma~6]{Yang}, with
divisor given by \cite[Lemma~7]{Yang}, \cite[Lemma~25]{GY}. For the pipeline's forms $c_0(0)=0$.

\subsection*{CM points}
Let $d<0$ be a discriminant, $d=\ell^2d_0$ with $d_0$ fundamental and $\ell\ge1$ the conductor,
$k=\Q(\sqrt d)$, $R_d=\Z+\ell\mathcal O_k$ the order of discriminant $d$, $h_0$, $\mu_0$ the class
number and number of roots of unity of $k$, and $\chi=\chi_{d_0}$ the Kronecker character of $k$,
primitive of conductor $|d_0|$. A CM point of discriminant $d$ on $X_0^D(N)$ is the fixed point
$\tau_d$ of $\phi(k^\times)$ for an optimal embedding $\phi:R_d\hookrightarrow\mathcal O$; $\CM(d)$
is the set of their images in $X$, and we assume $\CM(d)\neq\varnothing$. Put
\[
  \lambda=\phi(\sqrt d)\in L,\quad Q(\lambda)=|d|,\qquad U=\lambda^\perp\subset V,\qquad
  \Lp=L\cap\Q\lambda=\Z\lambda_0,\qquad \Lm=L\cap U,
\]
with $\lambda_0=\lambda$ or $\lambda/2$ according as $d\equiv1$ or $0\bmod4$
\parentref{the proof of \cite[Lemma~18]{Yang} for $\ell=1$; Lemma~\ref{lem:lambda0} in general}.
Then $\Lp\oplus\Lm\subset L\subset\Ldual\subset\Lp^\vee\oplus\Lm^\vee$, every $\eta\in\Ldual$
decomposes as $\eta=\eta_++\eta_-$ with $\eta_\pm\in L_\pm^\vee$, and $U$ is the negative definite
plane
\begin{equation}\label{eq:U}
  U\cong(k,\,-c\,N_{k/\Q})\quad\text{for some }c\in\Q_{>0}
\end{equation}
(Lemma~\ref{lem:split}), so that for every prime $p$: $U_p$ is isotropic iff $p$ splits in $k$, and
the quadratic character of $U$ is $\chi_{d_0}$. None of this depends on the conductor.

\subsection*{The Eisenstein series}
For $\nu\in\Lm^\vee/\Lm$ let $\varphi_\nu\in S(U(\A_f))$ be the characteristic function of
$\nu+\widehat{\Lm}$, and let $E(\tau,s;\varphi_\nu,1)=\sum_mA_\nu(s,m,v)q^m$ be the incoherent
Eisenstein series of weight $1$ attached to $(U,\varphi_\nu)$ as in \cite[\S4]{Yang},
\cite[\S2]{Schofer}: its finite components are those of $(U,Q)$ and its archimedean component
that of the positive definite plane $(U,-Q)$ \cite[after (2.18)]{BY}, so it vanishes identically
at $s=0$ and
\[
  A_\nu(s,m,v)=b_\nu(m,v)\,s+O(s^2).
\]
Following \cite[(11)]{Yang} we set
\begin{equation}\label{eq:kappaminus}
  \kappa^-_\nu(m)=\begin{cases}\lim_{v\to\infty}b_\nu(m,v)&m>0,\\
  \lim_{v\to\infty}\bigl(b_0(0,v)-\log v\bigr)&m=0,\ \nu=0,\\
  \lim_{v\to\infty}b_\nu(0,v)&m=0,\ \nu\neq0,\\ 0&m<0,\end{cases}
\end{equation}
where the third line is the one place we depart from \cite{Yang} and \cite[Lemma~2.21]{Schofer},
which set it to $0$; Theorem~\ref{thm:kappa}(c) computes it, and it is not $0$ in general. The
limits exist: for $m\neq0$ the $v$-dependence of $b_\nu(m,v)$ is exponentially small, and for
$m=0$, $\nu\neq0$ there is none at all (Lemma~\ref{lem:bconst}).

\section{The statement}\label{sec:statement}

\begin{definition}[the coefficient at a coset of the plane]\label{def:C}
For $\nu\in\Lm^\vee/\Lm$ and $n\in\Q$ put
\[
  C_\nu(n)=\sum_{\substack{x\in\Lp^\vee\\ x+\nu\in\Ldual}}c_{[x+\nu]}\bigl(-n-Q(x)\bigr),
\]
where $[\,\cdot\,]$ is the class in $\Ldual/L$ and the sum is over all $x$ in the fibre, both signs
included. It is finite because $\Lp$ is positive definite and $F$ has finitely many negative
exponents; it is the coefficient of $q^{-n}e_\nu$ in the contraction $\langle F,\theta_{\Lp}\rangle$
of $F$ against the theta function of $\Lp$ \cite[Cor.~3.2]{Schofer}.
\end{definition}

Only finitely many pairs $(\nu,n)$ with $n\ge0$ have $C_\nu(n)\neq0$, and $C_\nu(n)=0$ unless
$n\in Q(\nu)+\Z$ (as $Q(x+\nu)=Q(x)+Q(\nu)$ and $c_\eta(-m)=0$ unless $m\in Q(\eta)+\Z$).

\begin{theorem}[the CM value formula]\label{thm:main}
Let $F$ be as in Section~\ref{sec:setting} and $d<0$ any discriminant with $\CM(d)\ne\varnothing$
such that $\CM(d)$ does not meet the support of $\operatorname{div}\psi_F$. Then
\begin{equation}\label{eq:main}
  \sum_{\tau\in\CM(d)}\log\bigl|\psi_F(\tau)\,(\operatorname{Im}\tau)^{c_0(0)/2}\bigr|
  =-\frac{|\CM(d)|}{4}\Bigl(\sum_{\nu\in\Lm^\vee/\Lm}\ \sum_{n\ge0}C_\nu(n)\,\kappa^-_\nu(n)
     \;+\;c_0(0)\bigl(\Gamma'(1)+\log2\pi\bigr)\Bigr).
\end{equation}
Moreover $C_0(0)=c_0(0)+2\ord_{\tau_d}\Psi_F=c_0(0)$, so the term with $\kappa^-_0(0)$ is absent
whenever $c_0(0)=0$.
\end{theorem}

The right side of \eqref{eq:main} is the same number as Schofer's and Yang's
$\sum_{\eta\in\Ldual/L}\sum_{m\ge0}c_\eta(-m)\kappa_\eta(m)$ with $\kappa_\eta(m)$ of
\cite[(12)]{Yang} (Lemma~\ref{lem:regroup}), \emph{provided} $\kappa^-_\nu(0)$ at $\nu\neq0$ is
read as in \eqref{eq:kappaminus} wherever (12) invokes it. Writing it over the plane's cosets is
what makes the $m=0$ bookkeeping transparent: $\kappa^-_\nu(0)$ enters once, with the coefficient
$C_\nu(0)$, whatever the index $m$ of the term of $F$ it came from.

\begin{theorem}[the coefficients]\label{thm:kappa}
Let $\nu\in\Lm^\vee/\Lm$, $n\ge0$, $X_p=p^{-s}$, and let $W_{n,p}(s,\varphi_\nu)$ be the local
Whittaker function of $(U_p,\varphi_{\nu,p})$ in the normalisation of \cite[(4.4)]{KY}, so that
$W_{n,p}(s,\varphi_\nu)=\gamma_p\,|\det S|_p^{1/2}\,\mathcal W_{n,p}(X_p)$ with $\mathcal W_{n,p}$
the polynomial, or rational function, of \cite[Thms.~4.3, 4.4]{KY} computed from a Jordan
decomposition of $L_{-,p}$, $S$ the Gram matrix of $\langle\,,\rangle$ on $\Lm$ and $\prod_p\gamma_p=1$.
Let $\Sigma=\{p:\ L_{-,p}\text{ is not unimodular}\}$, which contains the primes of $d_0$, of
$\ell$ and of $DN$ that enter, and $\Sigma_n=\Sigma\cup\{p:\ p\mid n\}$ for $n>0$.
\begin{enumerate}
\item[(a)] ($n>0$.) $\kappa^-_\nu(n)=0$ unless $n\in Q(\nu)+\Z$. If $n\in Q(\nu)+\Z$, the number
of $p\in\Sigma_n$ with $W_{n,p}(0,\varphi_\nu)=0$ is at least one; if it is at least two,
$\kappa^-_\nu(n)=0$; if it is exactly one, $p'$ say, then
\[
  \kappa^-_\nu(n)=\frac{\mu_0\sqrt{|d_0|}}{h_0}\cdot\frac1{\sqrt{|\det S|}}\cdot
  \frac{\mathcal W'_{n,p'}(1)\prod_{p\in\Sigma_n,\,p\ne p'}\mathcal W_{n,p}(1)}
       {\prod_{p\in\Sigma_n}\bigl(1-\chi_{d_0}(p)/p\bigr)}\cdot\log p' ,
\]
with $\mathcal W'$ the derivative in $X_{p'}$; equivalently \cite[Lemma~16]{Yang} with $d_0$ in
place of $d$.
\item[(b)] ($n=0$, $\nu=0$.) With $G(s)=\prod_{p\in\Sigma}\dfrac{(1-\chi_{d_0}(p)p^{-s})\,
W_{0,p}(s,\varphi_0)/\gamma_p}{1-\chi_{d_0}(p)p^{-1-s}}$ and $\Lambda(s,\chi)=
(\pi/|d_0|)^{-(1+s)/2}\Gamma(\tfrac{1+s}2)L(s,\chi)$,
\[
  \kappa^-_0(0)=2\frac{\Lambda'(1,\chi_{d_0})}{\Lambda(1,\chi_{d_0})}-\frac{d}{ds}\log G(s)\Big|_{s=0},
\]
which is \cite[Lemma~17]{Yang} with $d_0$ in place of $d$.
\item[(c)] ($n=0$, $\nu\ne0$.) $\kappa^-_\nu(0)=0$ unless $Q(\nu)\in\Z$. Let
$\supp\nu=\{p:\ \nu_p\neq0\text{ in }L_{-,p}^\vee/L_{-,p}\}$. If $|\supp\nu|\ge2$ then
$\kappa^-_\nu(0)=0$. If $\supp\nu=\{p\}$ then the ratio $W_{0,p}(s,\varphi_\nu)/W_{0,p}(s,\varphi_0)$
is a rational function of $X_p$ vanishing at $X_p=1$, and
\[
  \kappa^-_\nu(0)=\log p\cdot\frac{d}{dX_p}\Bigl[\frac{W_{0,p}(s,\varphi_\nu)}{W_{0,p}(s,\varphi_0)}
  \Bigr]_{X_p=1}.
\]
\end{enumerate}
\end{theorem}

Everything on the right of (a)--(c) is a local computation on the actual lattice $L_{-,p}$; no
hypothesis on $\ell$, on $N$ beyond squarefreeness, or on the behaviour of the primes of $DN$ in
$k$ is used. The quantities $h_0$, $\mu_0$, $\chi_{d_0}$, $\Lambda(s,\chi_{d_0})$ are those of the
\emph{field}, for every conductor. Part~II evaluates the local factors in closed form for every
prime and coset that occurs, which is how the pipeline computes (c) (Proposition~\ref{prop:table}).

\begin{corollary}[the formula as implemented]\label{cor:impl}
Assume $c_0(0)=0$ and $\CM(d)$ off the divisor. Then
\begin{multline}\label{eq:impl}
  \sum_{\tau\in\CM(d)}\log|\psi_F(\tau)|=-\frac{|\CM(d)|}{4}\Bigl(
   \sum_{m>0}c_\infty(-m)\,\kappa_0(m)
  +\sum_{n>0}b_n\sum_{\substack{\eta\in\Ldual/L\\ Q(\eta)\equiv n/M\ (\Z)}}\kappa_\eta(n/M)
  \\
  +\sum_{\substack{\nu\in\Lm^\vee/\Lm,\ \nu\ne0\\ Q(\nu)\in\Z,\ |\supp\nu|=1}}
     \kappa^-_\nu(0)\,C_\nu(0)\Bigr),
\end{multline}
where in the first two sums $\kappa_\eta(m)=\sum_{\mu\in L/(\Lp+\Lm)}\sum_{x\in\eta_++\mu_++\Lp,
\ Q(x)<m}\kappa^-_{\eta_-+\mu_-}(m-Q(x))$ is \cite[(12)]{Yang} \emph{restricted to the pairs with
$Q(x)<m$}, and the pairs with $Q(x)=m$ are all collected in the last sum. The first sum is
\texttt{SchoferFormula} on the expansion at $\infty$, the second \texttt{SchoferFormula0} on the
expansion at the cusp $0$ (the bucketing of $\eta$ by $Q(\eta)\bmod\Z$ is \eqref{eq:GY24}), and the
last is the ``$m=0$ term'': $\kappa^-_\nu(0)$ from Proposition~\ref{prop:table} and $C_\nu(0)$ from
the fibre over $\nu$ (\texttt{M0FibreCorrection} at a prime of the conductor,
$\tfrac12c_\eta(0)-\sum_{k\ge1}c_\infty(-k^2Q(\lambda_0))$ at a level prime not dividing $\ell$,
Proposition~\ref{prop:levelprime}).
\end{corollary}

\begin{remark}[what the hypotheses exclude]
(1) A point of $\CM(d)$ on the divisor: there $C_0(0)=2\ord_{\tau_d}\Psi_F\ne0$ and the value is
infinite; what the formula can still give is the regularised value, in which each singular pair
$x$ with $Q(x)=m$, $c_{[x]}(-m)\neq0$ at the zero coset contributes $-c_{[x]}(-m)\log m$ plus a
constant independent of $m$ that cancels in every quotient of forms with the same order
(Remark~\ref{rem:divisor}). (2) Non-squarefree $N$: Theorems~\ref{thm:main} and \ref{thm:kappa}
hold verbatim --- nothing in Part~I uses squarefreeness --- but the local lattices of Part~II are
different, and the pipeline refuses such a level. (3) Weight $c_0(0)\ne0$: covered by
\eqref{eq:main}, with the extra terms of \cite[Thm.~B]{Yang}; the pipeline does not need them.
\end{remark}

\section{Proof of Theorem~\ref{thm:main}}\label{sec:proofmain}

The theorem is Schofer's \cite[Thm.~3.4(ii), Cor.~3.5(i)]{Schofer} for the lattice $L$, the
plane $U$ and the lattice $\Lm$, specialised to the Shimura curve as in \cite[Thm.~B]{Yang}, with
two points made explicit: the $m=0$ terms of his $(0,2)$ theorem are kept, instead of being
removed by \cite[Lemma~2.21]{Schofer}; and the normalisation of the CM cycle is checked for a
non-maximal order. We go through the chain.

\subsection*{Step 1: the $(1,2)$ theorem reduces to the $(0,2)$ theorem}
Let $\theta_{\Lp}(\tau)=\sum_{x\in\Lp^\vee}q^{Q(x)}e_{[x]}$ be the theta function of the positive
definite line $\Lp$, a holomorphic form of weight $\tfrac12$ and type $\rho_{\Lp}$. Schofer's
\cite[Prop.~3.1, Cor.~3.2]{Schofer}: the integral of the regularised theta lift of $F$ over the
CM cycle $Z(U)$ equals the regularised theta lift, for the lattice $\Lm$ of signature $(0,2)$,
of the contraction $\langle F,\theta_{\Lp}\rangle\in M^!_{1,\rho_{\Lm}}$, whose Fourier coefficients
are
\[
  \langle F,\theta_{\Lp}\rangle=\sum_{\nu\in\Lm^\vee/\Lm}\ \sum_{n}C_\nu(n)\,q^{-n}\,e_\nu,
  \qquad C_\nu(n)\text{ as in Definition~\ref{def:C}}.
\]
This uses only $L\subset\Lp^\vee\oplus\Lm^\vee$ and the compatibility of the Weil representations
of $L$ and of $\Lp\oplus\Lm$ under the inclusion $\Ldual/L\hookrightarrow(\Lp^\vee\oplus\Lm^\vee)/
(\Lp\oplus\Lm)$ \cite[Lemma~3.3]{Schofer}; no direct-sum hypothesis on $L$ itself. In
Schofer's formula the contraction is written as a sum over $\mu\in L/(\Lp\oplus\Lm)$ and over
$\eta$, which is the shape of \cite[(12)]{Yang}; Lemma~\ref{lem:regroup} shows the two
bookkeepings agree.

\subsection*{Step 2: the $(0,2)$ theorem}
For $G=\sum_{\nu,n}C_\nu(n)q^{-n}e_\nu$ weakly holomorphic of weight $1$ and type $\rho_{\Lm}$,
\cite[Thm.~3.4(ii)]{Schofer} evaluates the regularised integral $\Phi(z_0,h;G)$ over
$T(\Q)\backslash T(\A_f)/K_T$ as
\begin{equation}\label{eq:02}
  \int\Phi(z_0,h;G)\,dh=\lim_{t\to\infty}\Bigl[\,2\sum_{\nu}\sum_{n\ge0}C_\nu(n)\,b_\nu(n,t)
  \;-\;2\,C_0(0)\log t\Bigr],
\end{equation}
where $t$ is the truncation parameter of the regularisation and $b_\nu(n,t)$ the $s$-linear
coefficient of the Fourier coefficient $A_\nu(s,n,t)$ of the incoherent series. This is the
identity obtained by Stokes' theorem and the Siegel--Weil formula in \cite[(33) and the following
page]{Schofer}; at this point no coefficient has been discarded. In \cite{Schofer} the $n=0$ part
is then simplified by Lemma~2.21, ``$b_\nu(0,t)=0$ for $\nu\neq0$'', whose proof cites a local
vanishing valid only for unimodular lattices (Schofer assumes unimodularity from his \S4 on and says
so); we do not simplify, and instead record what the limit is.

\begin{lemma}\label{lem:bconst}
For $\nu\neq0$ the function $b_\nu(0,t)$ is independent of $t$; for $\nu=0$, $b_0(0,t)-\log t$ is
independent of $t$; for $n>0$, $b_\nu(n,t)$ converges as $t\to\infty$. Hence \eqref{eq:02} equals
$2\sum_{\nu}\sum_{n\ge0}C_\nu(n)\kappa^-_\nu(n)$ with $\kappa^-$ as in \eqref{eq:kappaminus}.
\end{lemma}
\begin{proof}
By \cite[Thm.~2.4]{KY} the constant term of $E(\tau,s;\varphi_\nu,1)$ is
$A_\nu(s,0,t)=\varphi_\nu(0)\,t^{s/2}+W_{0,\infty}(\tau,s)\prod_pW_{0,p}(s,\varphi_{\nu,p})$, and by
\cite[Prop.~2.3]{KY} the archimedean factor at $m=0$ is $W_{0,\infty}(\tau,s)=t^{-s/2}c_\infty(s)$
with $c_\infty$ holomorphic and nonzero near $s=0$. For $\nu\neq0$, $\varphi_\nu(0)=0$, and the
incoherent series vanishes at $s=0$, so $\prod_pW_{0,p}(s,\varphi_\nu)=O(s)$; then
$A_\nu(s,0,t)=t^{-s/2}c_\infty(s)\prod_pW_{0,p}=s\,c_\infty(0)\bigl(\prod_pW_{0,p}\bigr)'(0)+O(s^2)$,
and $t$ does not appear in the coefficient of $s$. For $\nu=0$, $A_0(s,0,t)=t^{s/2}-t^{-s/2}
c_\infty(s)\prod_pW_{0,p}(s,\varphi_0)$ with $c_\infty(0)\prod_pW_{0,p}(0,\varphi_0)=1$ (the series
vanishes at $s=0$), so $b_0(0,t)=\log t+(\text{constant})$. For $n>0$ the non-constant Fourier
coefficients of the Eisenstein series are $W_{n,\infty}(\tau,s)\prod_pW_{n,p}(s)$ with
$W_{n,\infty}(\tau,s)=q^n\,\omega(s,4\pi nt)$ and $\omega(s,y)\to$ a constant as $y\to\infty$
exponentially fast \cite[Prop.~2.3]{KY}, \cite[Lemma~2.20]{Schofer}; the limit exists.
\end{proof}

\subsection*{Step 3: from the integral to the sum over $\CM(d)$}
\cite[Cor.~3.5(i)]{Schofer} converts \eqref{eq:02} into $\sum_{z\in Z(U)}\log\lVert\Psi_F(z)\rVert^2
=-\frac{2}{\mathrm{vol}(K_T)}\cdot\frac12\int\Phi\,dh$, with the Petersson norm
$\lVert\Psi_F(z)\rVert=|\Psi_F(z)|\,|y|^{c_0(0)/2}$, and the constant $\Gamma'(1)+\log2\pi$ arises
from the relation between $\log\lVert\Psi_F\rVert$ and the regularised lift \cite[Thm.~13.3]{Borcherds},
exactly as in \cite[Thm.~B, Remark~15(1)]{Yang}. What has to be checked for a non-maximal order
is the measure. In Schofer's argument the CM cycle is the orbit $Z(U)=T(\Q)\backslash T(\A_f)/K_T$
with $T=GSpin(U)\cong k^\times$ and $K_T$ the stabiliser of $\widehat{\Lm}$ in $T(\A_f)$, i.e.\
the image of $\widehat{R_d}^\times$; its points are the CM points of discriminant $d$, that is, the
optimal embeddings of $R_d$ modulo the group, and under the Tamagawa measure on
$T(\Q)\backslash T(\A)$, of total volume $2$, every point of the orbit carries the same volume
because $T$ is commutative: each point has the stabiliser $R_d^\times/\{\pm1\}$ in $T(\Q)$. So
the theta integral over $T(\Q)\backslash T(\A)$ is $\frac{2}{|Z(U)|}\sum_{z\in Z(U)}$, as for the
maximal order, and the Siegel--Weil formula gives the Eisenstein series of the quadratic space $U$
with the Schwartz function of the lattice $\Lm$; the order enters only through $|Z(U)|$ and through
$\varphi_\nu$. Passing to $X=X_0^D(N)/W_{D,N}$ and to $\log|\psi_F|$ as in \cite[Remark~15(1)]{Yang}
gives \eqref{eq:main} with the factor $-|\CM(d)|/4$: the sum over $Z(U)$ becomes the sum over
$\CM(d)$, each point of $X$ being hit the same number of times, and the normalisation is that of
\cite[Thm.~B]{Yang}, which we take as the reference (it is the one against which every published
value has been checked; see Remark~\ref{rem:norm}).

\subsection*{Step 4: the term $C_0(0)$}
By Definition~\ref{def:C}, $C_0(0)=\sum_{x\in\Lp^\vee\cap\Ldual}c_{[x]}(-Q(x))=c_0(0)+
\sum_{x\neq0}c_{[x]}(-Q(x))$. By \cite[Thm.~A(2)]{Yang} the order of $\Psi_F$ at $z_\lambda$ is
$\sum_{r>0,\,r\lambda\in\Ldual}c_{r\lambda}(-r^2\langle\lambda,\lambda\rangle/2)=\sum_{x\in\Ldual
\cap\Q_{>0}\lambda}c_{[x]}(-Q(x))$, and $c_{[-x]}=c_{[x]}$, so the sum over $x\ne0$ is twice the
order, which is $0$ at a point off the divisor. This proves the last assertion of the theorem; note
that it needs no hypothesis on the individual terms, only on their sum, so a pole of $f$ at
$\infty$ at a CM exponent cancelled by the cusp-$0$ side is allowed (and then
$\kappa^-_0(0)$, a transcendental, $d$-dependent number, still drops out).

\begin{lemma}[the two bookkeepings agree]\label{lem:regroup}
With $\kappa_\eta(m)$ defined by \cite[(12)]{Yang} and $\kappa^-_\nu(0)$ for $\nu\ne0$ read as in
\eqref{eq:kappaminus},
\[
  \sum_{\eta\in\Ldual/L}\sum_{m\ge0}c_\eta(-m)\,\kappa_\eta(m)=\sum_{\nu\in\Lm^\vee/\Lm}\sum_{n\ge0}
  C_\nu(n)\,\kappa^-_\nu(n).
\]
\end{lemma}
\begin{proof}
The pairs $(\eta,\mu)$ with $\eta\in\Ldual/L$, $\mu\in L/(\Lp+\Lm)$ run over
$y=\eta+\mu\in\Ldual/(\Lp\oplus\Lm)$ exactly once, and $c_\eta(-m)=c_{[y]}(-m)$. For fixed $y$ the
inner sum of (12) runs over $x\in y_++\Lp$ and contributes $\kappa^-_{[y_-]}(m-Q(x))$; summing
$c_{[y]}(-m)\kappa^-_{[y_-]}(m-Q(x))$ over $m\ge0$ and substituting $n=m-Q(x)$ (so $n\ge0$ iff
$m\ge Q(x)$, the terms with $n<0$ vanishing) gives, for each $\nu=[y_-]$, the sum over the fibre
$\{x\in\Lp^\vee:\ x+\nu\in\Ldual\}$ of $c_{[x+\nu]}(-n-Q(x))$, which is $C_\nu(n)$.
\end{proof}

\begin{remark}[the normalisation, and what fixes it]\label{rem:norm}
The constant in front of \eqref{eq:main} is Yang's; Schofer states his Corollary~1.2 for
$\log\lVert\Psi\rVert^2$ with $-\deg Z(U)/2$, and the two agree ($\deg Z(U)=2h_k$ against
$\mathrm{vol}(K_T)=2/h_k$ in his normalisation of the measure). Errthum's version \cite{Err} has
no $\tfrac14$. The pipeline uses \eqref{eq:main} as written, and the normalisation has been
confirmed by every outside value we know of: the Chowla--Selberg values of \cite[\S4]{Yang}
(Examples~21, 24), the 254 values of the Guo--Yang CM tables and Table~45 of their arXiv version,
and the values forced by the divisors of the forms (\texttt{tests/\_offline/GuoYang\_*.m},
\texttt{tests/SchoferIsometry.m}, \texttt{tests/M0PoleSum.m}).
\end{remark}

\section{Proof of Theorem~\ref{thm:kappa}}\label{sec:proofkappa}

The three parts are three evaluations of the $s$-linear term of $A_\nu(s,n,v)$, all from the
factorisation of the Fourier coefficients of the Eisenstein series into local Whittaker functions
\cite[Thm.~2.4 and (4.4)]{KY}: for $n\in Q(\nu)+\Z$,
\begin{equation}\label{eq:factor}
  A_\nu(s,n,v)\,q^n=\delta_{\nu,n}\,v^{s/2}+W_{n,\infty}(\tau,s)\prod_{p<\infty}W_{n,p}(s,\varphi_{\nu,p}),
\end{equation}
$\delta_{\nu,n}=1$ if $(\nu,n)=(0,0)$ and $0$ otherwise, and $A_\nu(s,n,v)=0$ if $n\notin Q(\nu)+\Z$
(every $x\in\nu_p+L_{-,p}$ has $Q(x)\equiv Q(\nu_p)\bmod\Z_p$). The local factors are those of the
quadratic space $U_p$ and the function $\varphi_{\nu,p}$, computed by \cite[Thms.~4.3, 4.4]{KY}
from a Jordan decomposition of the $\Z_p$-lattice $L_{-,p}$; those theorems are stated for an
arbitrary lattice, and the pipeline evaluates them on the actual $L_{-,p}$ (\texttt{Wpoly},
\texttt{Wpoly2} in \texttt{SchoferFormula.m}; the agreement of these polynomials with the
representation densities obtained by counting was checked at every conductor plane of
$X_0^{15}(2)$ and $X_0^{10}(3)$, Remark~\ref{rem:count}). At a prime $p\notin\Sigma_n$ the lattice
is unimodular and $p\nmid n$, and the factor is $W_{n,p}=\gamma_p(1-\chi_{d_0}(p)p^{-1-s})$
\cite[(4.4), Thm.~4.3]{KY}, with $\chi_{d_0}$ the character of $U_p$ by \eqref{eq:U}.

\subsection*{(a) $n>0$}
Multiplying the unimodular factors together, $\prod_{p\notin\Sigma_n}(1-\chi(p)p^{-1-s})=
L(s+1,\chi)^{-1}\prod_{p\in\Sigma_n}(1-\chi(p)p^{-1-s})^{-1}$, and with the archimedean factor of
\cite[Prop.~2.3]{KY} at $n>0$, whose limit as $v\to\infty$ is $-2\pi q^n$ in weight $1$ at $s=0$
(the value \cite{Yang} uses in (13)), one gets
\begin{equation}\label{eq:euler}
  \lim_{v\to\infty}A_\nu(s,n,v)=-\frac{2\pi}{L(s+1,\chi_{d_0})}\prod_{p\in\Sigma_n}
  \frac{W_{n,p}(s,\varphi_\nu)}{1-\chi_{d_0}(p)p^{-1-s}}\quad(1+O(s)),
\end{equation}
where the $O(s)$ comes from the $s$-dependence of the archimedean factor and is irrelevant below
because the right side vanishes at $s=0$. It vanishes because the series is incoherent, so some
factor $W_{n,p}(0,\varphi_\nu)$, $p\in\Sigma_n$, vanishes; the derivative at $s=0$ of a product with
two vanishing factors is $0$, and with exactly one, at $p'$, it is the product of the other factors
at $s=0$ times $\frac{d}{ds}W_{n,p'}(s,\varphi_\nu)|_{s=0}=-\log p'\cdot\gamma_{p'}|\det S|_{p'}^{1/2}
\mathcal W'_{n,p'}(1)$. Inserting $L(1,\chi_{d_0})=2\pi h_0/(\mu_0\sqrt{|d_0|})$ (the class number
formula of the field), $\prod_p\gamma_p=1$ and $\prod_{p\in\Sigma_n}|\det S|_p^{1/2}=|\det S|^{-1/2}$
(all primes of $\det S$ lie in $\Sigma$) gives the display of (a). This is \cite[Lemma~16]{Yang}
with the field's invariants, which is what the pipeline evaluates (\texttt{kappaminus}; the earlier
code had the order's class number $h(R_d)$ in place of $h_0$, wrong by $h(R_d)/h_0$ at conductor
$4$, see Remark~\ref{rem:mpos}).

\subsection*{(b) $n=0$, $\nu=0$}
From \eqref{eq:factor} with $(\nu,n)=(0,0)$ and the unimodular factors,
\[
  A_0(s,0,v)=v^{s/2}-v^{-s/2}\frac{\Lambda(s,\chi_{d_0})}{\Lambda(s+1,\chi_{d_0})}\,G(s),
\]
where the archimedean factor and the completed $L$-function are assembled as in \cite[(14)]{Yang}
(his computation uses only that the unimodular Euler factors are $(1-\chi p^{-1-s})/(1-\chi p^{-s})$
up to $\gamma_p$, which is \cite[Thm.~4.3]{KY} at $m=0$ on a unimodular plane, and
$\prod_p\gamma_p=1$). Since $A_0(s,0,v)\equiv0$ at $s=0$, $G(0)=1$, and differentiating gives
$b_0(0,v)=\log v+2\Lambda'(1)/\Lambda(1)-G'(0)/G(0)$ by the functional equation
$\Lambda(s)=\Lambda(1-s)$, which is (b).

\subsection*{(c) $n=0$, $\nu\neq0$}
If $Q(\nu)\notin\Z$ then $A_\nu(s,0,v)\equiv0$ and $\kappa^-_\nu(0)=0$. Otherwise let
$\Phi(s)=A_0(s,0,v)-v^{s/2}=-v^{-s/2}c_\infty(s)\prod_pW_{0,p}(s,\varphi_0)$, so that $\Phi(0)=-1$
by (b). For $p\notin\supp\nu$ one has $\varphi_{\nu,p}=\varphi_{0,p}$, hence
\begin{equation}\label{eq:ratio}
  A_\nu(s,0,v)=\Phi(s)\prod_{p\in\supp\nu}\frac{W_{0,p}(s,\varphi_\nu)}{W_{0,p}(s,\varphi_0)},
\end{equation}
an identity of meromorphic functions (for $\mathrm{Re}\,s$ large it is the factorisation
\eqref{eq:factor} divided and multiplied by the same factors; both sides continue). Each factor of
the product vanishes at $s=0$: applying \eqref{eq:ratio} to the coset $\nu^{(p)}$ supported at
$p$ alone with $\nu^{(p)}_p=\nu_p$ --- a coset of $\Lm^\vee/\Lm$, since $\Lm^\vee/\Lm$ is the
product of its $p$-parts --- gives $A_{\nu^{(p)}}(s,0,v)=\Phi(s)\,W_{0,p}(s,\varphi_\nu)/
W_{0,p}(s,\varphi_0)$, which vanishes at $s=0$ by incoherence while $\Phi(0)\neq0$. So the product
in \eqref{eq:ratio} has a zero of order $\ge|\supp\nu|$ at $s=0$, which is the vanishing for
$|\supp\nu|\ge2$; for $\supp\nu=\{p\}$ the $s$-derivative at $0$ is $\Phi(0)$ times the derivative
of the ratio, and $dX_p/ds=-X_p\log p$ gives the display. The ratio is a rational function of
$X_p$ because both Whittaker functions are, by \cite[Thms.~4.3, 4.4]{KY}; when
$W_{0,p}(s,\varphi_0)$ has a pole at $X_p=1$ (the isotropic case, Lemma~\ref{lem:pole}) the ratio
has a zero there of the corresponding order and the formula reads as the limit. The normalising
constants $\gamma_p|\det S|_p^{1/2}$ cancel in the ratio, so any fixed normalisation of the local
factor may be used --- in Part~II the counting normalisation $W=(1-X)\sum_k\alpha_kX^k$ with
$\alpha_k=p^{-k}\#\{x\in(\nu+L_{-,p})/p^kL_{-,p}:\ Q(x)\equiv0\bmod p^k\}$.

Two facts about the Eisenstein series were used for the plane of a non-maximal order and for
$\varphi_\nu$ a nonzero coset: holomorphy at $s=0$ and vanishing there. Both are properties of the
quadratic space $U$ and of the Schwartz function, not of a lattice: \cite[Cor.~2.5]{KY} is stated
for an arbitrary $\varphi\in S(U(\A_f))$, and the series is incoherent because its finite
components come from $(U,Q)$ and its archimedean one from $(U,-Q)$, whatever lattice $\varphi$ is
the characteristic function of \cite[after (2.18)]{BY}; Bruinier--Ehlen--Yang run the same
argument for binary lattices of arbitrary discriminant, the orders $\mathcal O_D$ included
\cite[\S3]{BEY}. This completes the proof of Theorem~\ref{thm:kappa}. \qed

\begin{remark}[where the earlier notes sit in this proof]
Proposition~prop:kappa0 of \parentref{\texttt{kappa0-proof-standalone}} is (c) at a level prime
$p\nmid d_0$, $p\nmid\ell$, where the ratio is $(1-X)/(1+(p-2)X)$ (Lemma~\ref{lem:WN}) and the
derivative is $-1/(p-1)$; Proposition~prop:composite is the statement $|\supp\nu|\ge2\Rightarrow0$
together with that value prime by prime; Proposition~prop:fibre is (c) in general together with
Step~1 of Section~\ref{sec:proofmain}. The condition ``$N$ prime'' never entered the argument, only
the enumeration of the cosets; ``$d$ fundamental'' entered only through the local lattices, which
Part~II now describes at every conductor.
\end{remark}

\section{The $m=0$ terms made explicit}\label{sec:m0}

Theorem~\ref{thm:kappa}(c) and Definition~\ref{def:C} say what the last sum of \eqref{eq:impl} is;
this section says which pairs $(x,\nu)$ occur, so that the sum can be written out. Throughout
$\nu\neq0$ is a coset of $\Lm^\vee/\Lm$ with $Q(\nu)\in\Z$ and $\supp\nu=\{p\}$.

\begin{lemma}[the CM line]\label{lem:lambda0}
$\Lp=\Z\lambda_0$ with $Q(\lambda_0)=|d|$ if $d\equiv1\bmod4$ and $|d|/4$ if $d\equiv0\bmod4$;
$\Lp^\vee=\Z\lambda_0/(2Q(\lambda_0))$; and $\Ldual\cap\Q\lambda=\Z\lambda_0/g$ where $g$ is the
content of the linear form $\langle\lambda_0,\cdot\rangle$ on $L$, a divisor of $2Q(\lambda_0)$.
The vectors of the fibre over any $\nu$ are of the form $x=t\lambda_0$ with $t\in\frac1{2Q(\lambda_0)}\Z$,
and $Q(x)=t^2Q(\lambda_0)$; a vector $x=\lambda_0/p^j$ is the primitive CM vector of a point of
discriminant $d/p^{2j}$.
\end{lemma}
\begin{proof}
$\lambda=\phi(\sqrt d)$ and $R_d=\Z[(d+\sqrt d)/2]$, so $\lambda/2\in\mathcal O$ iff $d\equiv0\bmod4$,
and $\lambda/n\notin\mathcal O$ for $n>2$ by optimality (for $d\equiv0\bmod4$, $\lambda/(2n)\in
\mathcal O$ would put $\sqrt d/(2n)$, hence the order of discriminant $d/n^2$, into $\phi^{-1}(\mathcal O)$).
The rest is the definition of the dual and of $g$.
\end{proof}

\begin{proposition}[a level prime not dividing the conductor]\label{prop:levelprime}
Let $p\mid N$, $p\nmid\ell$. If $p\mid d_0$ there is no nonzero $\nu$ supported at $p$ with
$Q(\nu)\in\Z$ (Lemma~\ref{lem:Wp}, Remark~\ref{rem:ramified}): the prime contributes nothing. If
$p\nmid d_0$ then $p$ splits in $k$, $L_{-,p}$ is the $p$-scaled hyperbolic plane and
$L_p=L_{+,p}\oplus L_{-,p}$ (Lemma~\ref{lem:LN}), the $\nu$ in question are the $2(p-1)$ nonzero
isotropic vectors of $L_{-,p}^\vee/L_{-,p}\cong(\Z/p)^2$, each lies in $\Ldual$, each has
$\kappa^-_\nu(0)=-\log p/(p-1)$ (Lemma~\ref{lem:WN}), and
\[
  C_\nu(0)=c_{[\nu]}(0)+\sum_{x\in\Ldual\cap\Q\lambda,\,x\ne0}c_{[x+\nu]}(-Q(x))
         =c_{\eta_p}(0)-2\sum_{k\ge1}c_\infty\bigl(-k^2Q(\lambda_0)\bigr),
\]
where $c_{\eta_p}(0)$ is the common constant term of $F$ at the nonzero isotropic cosets of
$\Ldual/L$ of order $p$ (Lemma~\ref{lem:iso}). So the prime $p$ contributes to \eqref{eq:impl}
\[
  -\frac{|\CM(d)|}{4}\cdot2(p-1)\cdot\Bigl(-\frac{\log p}{p-1}\Bigr)C_\nu(0)
  =|\CM(d)|\,\log p\,\Bigl(\tfrac12c_{\eta_p}(0)-\sum_{k\ge1}c_\infty\bigl(-k^2Q(\lambda_0)\bigr)\Bigr).
\]
\end{proposition}
\begin{proof}
The structure statements are Part~II. For $C_\nu(0)$: since $\nu\in\Ldual$ (at $p$ by the direct
sum, elsewhere $\nu_q\in L_{-,q}\subset L_q$), the condition $x+\nu\in\Ldual$ is $x\in\Ldual\cap
\Q\lambda$, independent of $\nu$. The $x=0$ term is $c_{[\nu]}(0)$. For $x\neq0$: if $x\notin L$,
then $[x+\nu]$ and $[x]$ are nonzero classes with the same norm, so $c_{[x+\nu]}(-Q(x))=
c_{[x]}(-Q(x))$ by \eqref{eq:GY24}; if $x\in L$, $[x+\nu]=[\nu]\ne0$ while $[x]=0$, and
$c_{[\nu]}(-Q(x))=c_{[x]}(-Q(x))-c_\infty(-Q(x))$ again by \eqref{eq:GY24}. Hence
$\sum_{x\ne0}c_{[x+\nu]}(-Q(x))=\sum_{x\neq0}c_{[x]}(-Q(x))-\sum_{x\in L\cap\Q\lambda,\,x\ne0}
c_\infty(-Q(x))$; the first sum is $2\ord_{\tau_d}\Psi_F=0$ (Step~4 of Section~\ref{sec:proofmain})
and the second is $2\sum_{k\ge1}c_\infty(-k^2Q(\lambda_0))$, as $L\cap\Q\lambda=\Z\lambda_0$.
\end{proof}

The pole sum $\sum_kc_\infty(-k^2Q(\lambda_0))$ is nonzero exactly when $f$ has a pole at $\infty$
at a CM exponent of $d$ whose contribution to the order at $\tau_d$ is cancelled by the cusp-$0$
side. It vanishes identically when $f$ is holomorphic at the cusp $0$ (even $D$), and it has been
zero at every evaluated point of the model set (\texttt{M0InfinityPoleSum}); it is nevertheless
part of the formula.

\begin{proposition}[a prime of the conductor]\label{prop:conductor}
Let $p\mid\ell$, $p^k\,\|\,\ell$. Then $L_p\neq L_{+,p}\oplus L_{-,p}$, the cosets $\nu$ in
question are the $p^k-1$ nonzero cosets of integral norm listed in Lemmas~\ref{lem:conductor},
\ref{lem:conductor2}, \ref{lem:unimod}, \ref{lem:unimod2} and \ref{lem:ramlevel} ($\nu_r$ of
order $p^{k-\rho}$, $0\le\rho<k$), their $\kappa^-_\nu(0)$ are the closed forms of
Lemma~\ref{lem:Wcond} and of those lemmas, and the fibre over $\nu_r$ consists of the vectors
$x=t\lambda_0$ with $x+\nu_r\in\Ldual$, which are: $x=0$ when $\nu_r\in\Ldual$ (only the cosets of
order $p$ at a level prime, Lemma~\ref{lem:conductor}(4), \ref{lem:conductor2}(4),
\ref{lem:ramlevel}; none at a prime outside the level, Lemma~\ref{lem:unimod}(2),
\ref{lem:unimod2}(2)), contributing $c_{[\nu_r]}(0)=c_{\eta_p}(0)$; and the vectors
$x=\pm j\lambda_0/p^{k-\rho}$, $p\nmid j$, with the congruence on $j$ given in those lemmas,
contributing the cusp-$0$ coefficient $b_{MQ(x)}$ of $f$ at the exponent $Q(x)$ when $[x+\nu_r]\ne0$
and the full zero-coset coefficient $c_0(-Q(x))=c_\infty(-Q(x))+b_{MQ(x)}$ when $x+\nu_r\in L$.
The contribution of $p$ to \eqref{eq:impl} is $-\frac{|\CM(d)|}4\sum_{r\ne0}\kappa^-_{\nu_r}(0)
\,C_{\nu_r}(0)$ with these ingredients; this is what \texttt{M0FibreCorrection} evaluates, by
enumerating the $p$-primary integral cosets of $\Lm^\vee/\Lm$ and the fibre over each from the
lattice.
\end{proposition}
\begin{proof}
Theorem~\ref{thm:kappa}(c), Definition~\ref{def:C} and Part~II; the coefficient at a nonzero class
is \eqref{eq:GY24}.
\end{proof}

The vectors $x=\lambda_0/p^j$ of the fibre are the primitive CM vectors of the discriminants
$d/p^{2j}$ (Lemma~\ref{lem:lambda0}): at a conductor prime the formula sees the poles of $f$ at the
CM points of the smaller discriminants $d/p^2,d/p^4,\dots$, which no $d$-independent multiplier can
reproduce, and which is why the earlier rule ``add $\tfrac12c_\eta(0)\log p$ when $p$ splits''
failed at conductor $4$ (Remark~\ref{rem:evidence}).

\begin{proposition}[the table of local factors]\label{prop:table}
For $N$ squarefree prime to $D$ and any $d$ with $\CM(d)\neq\varnothing$, the primes $p$ at which
a nonzero $\nu$ with $Q(\nu)\in\Z$ and $\supp\nu=\{p\}$ exists, and the values of $\kappa^-_\nu(0)$,
are the following ($\epsilon=\chi_{d_0}(p)$, $p^k\,\|\,\ell$, $\nu$ of order $p^{k-\rho}$ when $p\mid\ell$):
\begin{center}\small
\begin{tabular}{lll}
\toprule
prime & $\nu$ & $\kappa^-_\nu(0)/\log p$\\
\midrule
$p\nmid DN\ell$ & none & ---\\
$p\mid D$ & none (anisotropic plane, Lemma~\ref{lem:Wp}) & ---\\
$p\mid N$, $p\nmid\ell$, $p\mid d_0$ & none (Remark~\ref{rem:ramified}) & ---\\
$p\mid N$, $p\nmid\ell$, $p\nmid d_0$ & $2(p-1)$ cosets of order $p$ & $-1/(p-1)$\\
$p\mid N$, $p\mid\ell$, $\epsilon=1$ & $p^k-1$ cosets, $0\le\rho<k$ & $-2p^{\rho-k+1}/(p-1)$\\
$p\mid N$, $p\mid\ell$, $\epsilon=-1$ & $p^k-1$ cosets & $-\dfrac{2}{p^{k-1}(p+1)}\Bigl(p^\rho+\dfrac{2(p^\rho-1)}{p-1}\Bigr)$\\
$p\mid N$, $p\mid\ell$, $\epsilon=0$ & $p^k-1$ cosets (Lemma~\ref{lem:ramlevel}) & $-2(p^{\rho+1}-1)/((p-1)p^k)$\\
$p\nmid DN$, $p\mid\ell$, $\epsilon=\pm1$ & $p^k-1$ cosets, none in $\Ldual$ & as for $p\mid N$, $p\mid\ell$\\
$p\nmid DN$, $p\mid\ell$, $\epsilon=0$ & $p^k-1$ cosets, none in $\Ldual$ & $-2(p^{\rho+1}-1)/((p-1)p^k)$\\
\bottomrule
\end{tabular}
\end{center}
At $p=2$ the same values hold with the cosets of Lemmas~\ref{lem:conductor2} and \ref{lem:unimod2}.
\end{proposition}
\begin{proof}
Lemmas~\ref{lem:LN}, \ref{lem:WN}, \ref{lem:Wp}, \ref{lem:conductor}, \ref{lem:conductor2},
\ref{lem:Wcond}, \ref{lem:unimod}, \ref{lem:unimod2}, \ref{lem:ramlevel} and
Remark~\ref{rem:ramified}, each applied through Theorem~\ref{thm:kappa}(c). A prime of $D$ cannot
divide $\ell$ (at a prime of $D$ only the maximal order of $k_p$ embeds optimally into the maximal
order of the division algebra).
\end{proof}

\begin{remark}[the constant terms $c_{\eta}(0)$]\label{rem:ceta}
The $x=0$ terms need the constant terms of $F$ at the nonzero isotropic cosets of $\Ldual/L$,
which are not given by \eqref{eq:GY24}. They are constant on each support class (the orthogonal
group of the discriminant form acts transitively on the isotropic vectors of each $p$-part and
commutes with $\rho_L$, Lemma~\ref{lem:iso}), so one value per prime $p\mid N$ suffices, and they
are computed exactly, as elements of $\Q(\zeta_{24M})$ reduced to $\Q$, by
\texttt{M0MultipliersAlgebraic} (Remark~\ref{rem:exact}); the numerical coset evaluation
\texttt{M0MultipliersBySupport} is the cross-check. At a prime of $N$ the value differs from class
to class in general (on $X_0^6(35)$: $-\tfrac{23}4$ at $\{5\}$, $-\tfrac{39}4$ at $\{7\}$, $0$ at
$\{5,7\}$ for one form), which is why the multiplier is indexed by the prime.
\end{remark}

\begin{remark}[the counting normalisation and its check]\label{rem:count}
In Part~II the local factors at $m=0$ are written in the counting normalisation $W_{0,p}=(1-X)
\sum_k\alpha_kX^k$. Theorem~\ref{thm:kappa}(c) uses only ratios, so the normalisation is immaterial
there; for (a) the pipeline uses \cite[Thms.~4.3, 4.4]{KY} on $L_{-,p}$ with the factor
$|\det S|_p^{1/2}$, and these polynomials agree with the counting densities at every $(d,p,m)$
tried on the conductor planes of $X_0^{15}(2)$ ($p=2$; $d=-15,-60,-240,-28,-12,-48$, $m\le6$)
and at the unimodular prime $7$ (\texttt{level-p2/wnorm.m}). The Euler-product identity behind
(a) --- $L(1,\chi_{d_0})^{-1}$ times the correction $(1-\chi_{d_0}(p)/p)^{-1}$ at every prime of
$\Sigma_n$, the conductor primes included --- is also confirmed by Siegel's genus-average identity
on the definite plane, which holds with one universal constant at $d=-15,-60,-240,-48$ with these
corrections and fails without them (\texttt{level-p2/siegel.m}).
\end{remark}

\begin{remark}[the earlier error at non-fundamental $d$]\label{rem:mpos}
The pipeline once used $h(R_d)$, the class number of the order, in (a); relative to the correct
$h_0$ this is off by $h(R_d)/h_0$, invisible at conductor $2$ when $p$ splits ($h(-60)=h(-15)$,
$h(-28)=h(-7)$) and exactly halving every value at conductor $4$. It amounts to using
$L(1,\chi_d)$ of the imprimitive character in place of $L(1,\chi_{d_0})$, i.e.\ to applying the
correction at $p\mid\ell$ twice (the correction is already in the Euler product, since every prime
of $\ell$ divides $\det S$). Magma's \texttt{KroneckerCharacter(d)} is the primitive character of
conductor $|d_0|$, which is what (a) needs.
\end{remark}

\begin{remark}[evidence]\label{rem:evidence}
What the formula has been checked against, beyond the fundamental-discriminant tables: on
$X_0^{15}(2)$, the nine forms of the model set at the Table-45 points of conductors $1$, $2$ and
$4$ ($d=-7,-15,-28,-60,-240,-48$; $46$ values, each forced by the form's divisor and the
Hauptmodul), where the fibre sum of Proposition~\ref{prop:conductor} is exact and the earlier rule
was right on $8$ of $18$ conductor-$4$ values; at $d=-588=14^2\cdot(-3)$ ($7$ a conductor prime
outside the level), the class polynomial of the three star points is a rational cubic cutting out
the field of definition of the points only with the term of Lemma~\ref{lem:unimod}; the published
values $s(-27)$ on $X_0^{58}(1)$ and $X_0^{55}(1)$ ($3\mid d_0$, $3\mid\ell$), $s(-100)$ on
$X_0^{21}(2)$, $s(-16)$ on $X_0^{21}(2)$ ($2$ a level prime dividing $d_0$ and $\ell$,
Lemma~\ref{lem:ramlevel}), and $s(-28),s(-60),s(-12),s(-16),s(-315)$ on the odd-level bases
$X_0^{35}(1)$, $X_0^{39}(1)$, $X_0^{51}(1)$, $X_0^{55}(1)$, $X_0^{57}(1)$, all reproduced
(\texttt{tests/M0PoleSum.m}, \texttt{tests/\_offline/GuoYang\_*.m}); and at composite level, the
models of $X_0^6(35)$, $X_0^{10}(21)$, $X_0^{14}(15)$ built with Proposition~\ref{prop:levelprime}
prime by prime have the point counts of the trace formula at six good primes.
\end{remark}

\begin{remark}[a point on the divisor]\label{rem:divisor}
At $\tau_d$ on the divisor a single value is infinite but a quotient of two forms with the same
order there is finite. In the formula the point shows up as pairs $(x,0)$ with $Q(x)=m$ and
$c_{[x]}(-m)\ne0$, i.e.\ $C_0(0)\neq0$; the regularised Green function leaves at such a pair the
constant $-\log m$ per unit coefficient (Kudla's $\beta(2\pi vR)=-\gamma-\log(2\pi vR)+\dots$ with
$R$ proportional to $m$ at the point) plus a constant independent of $m$, which cancels in every
quotient. With that rule the pipeline reproduces $f_{-2}/f_{-1}=1/20$ and $f_{10}/f_{-1}^3=
9/(2^{16}\cdot5)$ at $\tau_{-12}$ on $X_0^{15}(2)$ (Table~45 and the divisors) and the $36$
quotients of the nine forms of $X_0^{21}(2)$ at $\tau_{-4}$ (\texttt{tests/DivisorPoint\_21\_2.m});
the two heuristics taken from Yang's code that it replaces gave $2^{2/3}/5$ for the first. Off the
divisor no such pair occurs at any evaluated point.
\end{remark}

\part{The local lattices and their $m=0$ factors}

What follows is the case-by-case evaluation of the local factors that Theorem~\ref{thm:kappa}(c)
and Proposition~\ref{prop:table} rest on, carried over from \parentref{\texttt{kappa0-proof-standalone.tex}}
with the labels of that note. Throughout, $\lambda=\phi(\sqrt d)$ for an optimal embedding of
$R_d$, $\langle\alpha,\beta\rangle=\tr(\alpha\beta')$, $Q=\langle\cdot,\cdot\rangle/2$, and the
counting normalisation of the $m=0$ local factors is $W_{0,p}(s,\varphi_\nu)=(1-X)\sum_{j\ge0}
\alpha_j(\nu)X^j$, $X=p^{-s}$, $\alpha_j(\nu)=p^{-j}\#\{x\in(\nu+L_{-,p})/p^jL_{-,p}:\ Q(x)\equiv
0\bmod p^j\}$. In this part ``$\kappa^-_\nu(0)$'' is always obtained from these series through
Theorem~\ref{thm:kappa}(c).

\section{Fundamental discriminant}
In this section the prime of the level under discussion is written $N$; at a squarefree level the statements hold prime by prime, the cosets of $\Lm^\vee/\Lm$ being the products of their $p$-parts. The discriminant $d$ is fundamental, $N\nmid d$ unless said otherwise, $\lambda=\phi(\sqrt d)$, $\langle\lambda,\lambda\rangle=2|d|$.

\begin{lemma}[the lattice at the level prime]\label{lem:LN}
$L_{-,N}$ is the $N$-scaled hyperbolic plane, $N\begin{pmatrix}0&1\\1&0\end{pmatrix}$ up to a
unit; $L_{-,N}^\vee/L_{-,N}\cong(\Z/N)^2$ with $Q(x,y)=xy/N\bmod\Z$, and the nonzero isotropic
cosets are the $2(N-1)$ nonzero vectors on the two lines $x=0$, $y=0$. Moreover
$L_N=L_{+,N}\oplus L_{-,N}$.
\end{lemma}
\begin{proof}
\cite[Lemma~18(1)]{Yang}: for odd $N\nmid d$ the Gram matrix of $L_{-,N}$ is
$N\,\mathrm{diag}(\epsilon_1,\epsilon_2)$ with $\epsilon_1\epsilon_2=-d$, and $\chi_d(N)=1$ makes
$-\epsilon_1\epsilon_2$ a square, so the plane is split. For $N=2$ (then $d\equiv1\bmod8$)
\cite[Lemma~18(3)]{Yang} gives $2\epsilon\begin{pmatrix}0&1\\1&0\end{pmatrix}$ directly (Yang
writes out the $2$-adic proof only for $2\nmid DN$; we checked this case: $L_{-,2}$ has basis
$\begin{pmatrix}-c_3/c_1&1\\0&c_3/c_1\end{pmatrix}$,
$\begin{pmatrix}-c_2/c_1&0\\2&c_2/c_1\end{pmatrix}$ in the notation below, with Gram matrix twice
an even form with odd off-diagonal entry and determinant $\equiv-1\bmod8$).

For the direct sum it suffices that $\langle\lambda,L_N\rangle\subset\langle\lambda,\lambda\rangle
\Z_N$: then the orthogonal projection $x\mapsto\frac{\langle x,\lambda\rangle}{\langle\lambda,
\lambda\rangle}\lambda$ maps $L_N$ into $\Z_N\lambda=L_{+,N}$ and $x-\mathrm{pr}(x)\in L_N\cap
\lambda^\perp=L_{-,N}$. For odd $N$, $2\lvert d\rvert$ is an $N$-unit. For $N=2$:
$\mathcal O\otimes\Z_2=\begin{pmatrix}\Z_2&\Z_2\\2\Z_2&\Z_2\end{pmatrix}$ (remark in the proof of
\cite[Lemma~18]{Yang}), $\lambda=\begin{pmatrix}c_1&c_2\\2c_3&-c_1\end{pmatrix}$, and optimality
of the embedding of $\mathcal O_d$, $d\equiv1\bmod8$, means $(1+\lambda)/2\in\mathcal O\otimes
\Z_2$, i.e.\ $c_1$ odd and $c_2,c_3$ even; for $\alpha=\begin{pmatrix}a&b\\2c&-a\end{pmatrix}\in
L_2$, $\langle\alpha,\lambda\rangle=-\tr(\alpha\lambda)=-2(ac_1+bc_3+cc_2)\in2\Z_2=
2\lvert d\rvert\Z_2$.
\end{proof}

Note that $L_{+,N}=\Z_N\lambda$ is unimodular for odd $N$ but not for $N=2$, where
$L_{+,2}^\vee/L_{+,2}\cong\Z/2$ is generated by $\lambda/2$ with $Q(\lambda/2)=\lvert d\rvert/4
\notin\Z_2$.

\begin{lemma}[the $m=0$ local factor at $N$]\label{lem:WN}
With $X=N^{-s}$ and $\gamma_N$ the local splitting index of \cite[(4.3)]{KY}, for $\nu\in
L_{-,N}^\vee/L_{-,N}$,
\[
  W_{0,N}(s,\varphi_0)=\frac{\gamma_N}{N}\cdot\frac{1+(N-2)X}{1-X},\qquad
  W_{0,N}(s,\varphi_\nu)=\frac{\gamma_N}{N}\ \text{($\nu\neq0$ isotropic)},\qquad
  W_{0,N}(s,\varphi_\nu)=0\ \text{($\nu\neq0$ anisotropic)}.
\]
\end{lemma}
\begin{proof}
The first is \cite[(15)]{Yang} with $\chi_d(N)=1$ for odd $N$, where $\frac{1-X/N}{N-X}=\frac1N$,
and \cite[(16)]{Yang} for $N=2$. For the others apply \cite[(4.4) and Thm.~4.3]{KY} ($N$ odd; KY's
matrix $S$ is the matrix of the bilinear form with $Q=S/2$, as is Yang's Gram matrix, so
$S=N\,\mathrm{diag}(\epsilon_1,\epsilon_2)$, $l_1=l_2=1$, $s_0=n/2-1=0$) to $\nu=(a/N,b/N)$ with
$m=0$: $Q(\nu)=(\epsilon_1a^2+\epsilon_2b^2)/(2N)$ lies in $\Z_N$ exactly when $\nu$ is isotropic
(for odd $N$ the $2$ is a unit and may be dropped), so
$W=0$ otherwise by Thm.~4.3(i); when it is, $-\epsilon_1\epsilon_2$ being a square forces $a,b$
both units, so $H_\nu=\varnothing$, $K_0(\nu)=\min(l_i+\ord_N\nu_i)=0$, and $a_\nu=\ord_N(t_\nu)
\ge0=K_0$ puts us in the second case of Thm.~4.3(ii), whose sum over $0<k\le K_0$ is empty:
$W_N(s,0,\nu)=1$. For $N=2$, \cite[Thm.~4.4]{KY} with the Gram matrix $2\epsilon\begin{pmatrix}
0&1\\1&0\end{pmatrix}$, i.e.\ KY's (4.9) with $H=N=0$, $M=1$, $m_1=1$: $\nu'=(a/2,b/2)\notin
\Z_2^2$ gives $M_\nu=\varnothing$, $K_0=m_1+\ord_2\nu'=0$, both sums in (ii) are empty, and
$W_2(s,0,\nu)=1$ when $Q(\nu)=ab/2\in\Z_2$, i.e.\ at the two nonzero isotropic cosets. In every
case the normalising constant of \cite[(4.4)]{KY} is $\gamma_N\lvert\det S\rvert_N^{1/2}=
\gamma_N/N$, the same for all $\nu$, so it cancels in the ratio used below. \textup{Lemma~lem:m0loc of the preprint}
is the same object in a counting normalisation (its unimodular calibration is Thm.~4.3 with
$l_i=0$), and its third and fourth rows are these values.
\end{proof}

\begin{lemma}[no other nonzero coset contributes]\label{lem:Wp}
For $p\neq N$ every nonzero $\nu_p\in L_{-,p}^\vee/L_{-,p}$ is anisotropic, $Q(\nu_p)\notin\Z_p$.
Consequently $W_{0,p}(s,\varphi_{\nu,p})=0$ for such $\nu_p$, and a nonzero isotropic coset of
$\Lm^\vee/\Lm$ is supported at $N$.
\end{lemma}
\begin{proof}
$W_{0,p}=0$ follows from anisotropy by \cite[Thm.~4.3(i)]{KY} ($p$ odd) and \cite[Thm.~4.4(i)]{KY}
($p=2$). Anisotropy is read off \cite[Lemma~18]{Yang}; $d$ fundamental gives $p\,\|\,d$ for odd
$p\mid d$.

$p$ odd, $p\nmid dDN$: unimodular. $p$ odd, $p\mid d$: Gram $\mathrm{diag}(\epsilon_1,\epsilon_2)$
with, say, $\epsilon_2=p\epsilon_2'$; $L^\vee/L\cong\Z/p$ generated by $(0,1/p)$, and
$Q(0,b/p)=\epsilon_2'b^2/(2p)\notin\Z_p$ for $p\nmid b$. $p$ odd, $p\mid D$, $p\nmid d$: Gram
$p\,\mathrm{diag}(\epsilon_1,\epsilon_2)$ with $\epsilon_1\epsilon_2=-d$; $p$ is inert in $k$, so
$d$ is a non-square mod $p$ (the Hilbert-symbol condition $(-\mu_1,-\mu_2p)_p=-1$ in Yang's
proof), and $Q(a/p,b/p)=(\epsilon_1a^2+\epsilon_2b^2)/(2p)\in\Z_p$ forces $a\equiv b\equiv0$.

$p=2\nmid d$. If $d\equiv1\bmod8$ then $2\nmid D$ and, as $2\neq N$, Lemma~18(3) gives
$\epsilon\begin{pmatrix}0&1\\1&0\end{pmatrix}$, unimodular. If $d\equiv5\bmod8$, Lemma~18(4) gives
$\epsilon\begin{pmatrix}2&1\\1&2\end{pmatrix}$ for $2\nmid D$, unimodular, and $2\epsilon
\begin{pmatrix}2&1\\1&2\end{pmatrix}$ for $2\mid D$, where $L^\vee/L\cong(\Z/2)^2$ and
$Q(a/2,b/2)=\epsilon(a^2+ab+b^2)/2$ with odd numerator for $(a,b)\not\equiv(0,0)$. (This case is
again one Yang only sketches; with $(1+\lambda)/2$ in the maximal order, $a,b,c$ are all odd in
$\lambda=aI+bJ+cIJ$ below, and the elimination below gives diagonal entries $\equiv2\bmod8$, a
unit off-diagonal and determinant $\equiv3\bmod8$.)

$p=2\mid d$: Lemma~18(2) gives $2\,\mathrm{diag}(\epsilon_1,\epsilon_2)$ with
$\epsilon_1\epsilon_2=-d/4$, $Q(x,y)=\epsilon_1x^2+\epsilon_2y^2$, and $d$ fundamental means
$d/4\equiv2,3\bmod4$. If $d\equiv12\bmod16$ both $\epsilon_i$ are units with $\epsilon_1\equiv
\epsilon_2\bmod4$, and the three nonzero cosets of $(\Z/2)^2$ have $Q=\epsilon_1/4$,
$\epsilon_2/4$, $(\epsilon_1+\epsilon_2)/4\equiv\epsilon_1/2$. If $d\equiv8\bmod16$, say
$\epsilon_2=2\epsilon'$; $L^\vee/L\cong\Z/2\times\Z/4$ and $Q(a/2,b/4)=\epsilon_1a^2/4+
\epsilon'b^2/8$ has $2$-adic valuation $-3$ for $b$ odd, $-2$ or $-1$ for $b\equiv2\bmod4$, and
equals $\epsilon_1/4$ for $b\equiv0$, $a$ odd. None is in $\Z_2$.

Yang writes out the proof of Lemma~18(2) only for $2\nmid DN$; the case $2\mid N$, $2\mid d$ does
not occur here, and the case $2\mid D$, $2\mid d$ is proved as follows. $B\otimes\Q_2=
\bigl(\frac{-1,-1}{\Q_2}\bigr)$ with maximal order $\Z_2+\Z_2I+\Z_2J+\Z_2(1+I+J+IJ)/2$, whose
trace-zero part is $L_2=\Z_2I+\Z_2J+\Z_2IJ$ with $Q=x^2+y^2+z^2$. Since $d\equiv0\bmod4$,
$\mathcal O_d=\Z[\sqrt d/2]$ and $\Lp=\Z\lambda_0$ with $\lambda_0=\lambda/2=aI+bJ+cIJ$,
$a^2+b^2+c^2=m:=-d/4\equiv1$ or $2\bmod4$, and $L_{-,2}=\{xI+yJ+zIJ:ax+by+cz=0\}$. If
$m\equiv1\bmod4$ exactly one of $a,b,c$ is odd, say $a$; eliminating $x=-(by+cz)/a$,
$Q|_{L_{-,2}}=u_1y^2+2w\,yz+u_2z^2$ with $u_1=b^2/a^2+1$, $u_2=c^2/a^2+1$ units and $w=bc/a^2\in
4\Z_2$; $y'=y+(w/u_1)z$ is an integral change of basis and $Q=u_1y'^2+(u_2-w^2/u_1)z^2$ with
$u_1u_2-w^2=m/a^2$. If $m\equiv2\bmod4$ two of $a,b,c$ are odd, say $a,b$; then $u_1\in2\Z_2^\times$,
$u_2$ is a unit, $w\in2\Z_2$, and completing the square on $z$ instead gives $Q=u_2z'^2+
(u_1-w^2/u_2)y^2$, product again $m/a^2$. Rescaling one basis vector by the unit $a$ makes the
product $-d/4$ exactly, which is Lemma~18(2).
\end{proof}

\begin{lemma}[isotropic cosets of $\Ldual/L$]\label{lem:iso}
For $N$ prime the isotropic cosets of $\Ldual/L$ are supported at $N$ and number $2N-1$: under
$L_N=L_{+,N}\oplus L_{-,N}$ they are the cosets $0\oplus\nu$ with $\nu$ an isotropic coset of
$L_{-,N}^\vee/L_{-,N}$. All $2N-2$ nonzero ones carry the same constant term $c_\eta(0)$ of $F_f$.
\end{lemma}
\begin{proof}
A coset is isotropic iff each $p$-part is, and $\Ldual/L$ does not depend on $d$, so the
computation is on $L_p$ itself. For $p$ odd, $p\nmid DN$, $L_p$ is unimodular. For $p$ odd,
$p\mid DN$, the proof of \cite[Lemma~18]{Yang} gives a basis of $L_p$ with Gram matrix
$\mathrm{diag}(2\mu_1,2\mu_2p,2\mu_1\mu_2p)$, $\mu_i$ units, so $(\Ldual/L)_p\cong(\Z/p)^2$ with
$Q(ae_2/p+be_3/p)=\mu_2(a^2+\mu_1b^2)/p$: anisotropic iff $-\mu_1$ is a non-square mod $p$, i.e.\
iff $(-\mu_1,-\mu_2p)_p=-1$, which Yang's proof identifies with $p\mid D$; for $p\mid N$ it is the
hyperbolic plane over $\mathbb F_p$, consistent with Lemma~\ref{lem:LN}. For $p=2\nmid DN$, $L_2$
is the trace-zero part of $M_2(\Z_2)$, with Gram $(-2)\oplus\begin{pmatrix}0&-1\\-1&0\end{pmatrix}$
in the basis $\mathrm{diag}(1,-1),e_{12},e_{21}$, so $(\Ldual/L)_2\cong\Z/2$ generated by
$\tfrac12\mathrm{diag}(1,-1)$ with $Q=-1/4$. For $p=2\mid D$, $L_2=\Z_2I+\Z_2J+\Z_2IJ$ as in the
proof of Lemma~\ref{lem:Wp}, $(\Ldual/L)_2\cong(\Z/2)^3$ with $Q=(a^2+b^2+c^2)/4\in\{\tfrac14,
\tfrac12,\tfrac34\}$ on nonzero elements. At $p=N$ odd, $L_{+,N}$ is unimodular and
$(\Ldual/L)_N=L_{-,N}^\vee/L_{-,N}$, with $2N-1$ isotropic vectors. At $N=2$, the orthogonal direct
sum gives $(\Ldual/L)_2=L_{+,2}^\vee/L_{+,2}\oplus L_{-,2}^\vee/L_{-,2}\cong\Z/2\oplus(\Z/2)^2$; the
elements $\lambda/2\oplus\nu$ have $Q=\lvert d\rvert/4+Q(\nu)$ with $Q(\nu)\in\{0,\tfrac12\}\bmod
\Z$ and $\lvert d\rvert$ odd, never integral, so the isotropic elements are the $0\oplus\nu$ with
$\nu$ isotropic: again $2N-1=3$.

For the constant terms: an isometry $\sigma$ of the discriminant form permutes the $e_\eta$, fixes
$e_0$, and commutes with $\rho_L(S)$ and $\rho_L(T)$ because it preserves $Q$ and the pairing,
hence with every $\rho_L(\tilde\gamma)$; so $\sigma F_f=F_f$ and $c_{\sigma\eta}(0)=c_\eta(0)$. The
orthogonal group of the hyperbolic plane over $\mathbb F_N$ contains the torus $(x,y)\mapsto
(ax,a^{-1}y)$ and the swap, and acts transitively on the $2(N-1)$ nonzero isotropic vectors; these
isometries of $(\Ldual/L)_N$ extend by the identity on the other $p$-parts.
\end{proof}

\begin{remark}[a prime of $N$ dividing $d$]\label{rem:ramified}
If $p\mid N$ and $p\mid d$, nothing is dropped at $p$: $L_{-,p}$ is anisotropic, so no nonzero
isotropic coset is supported there and $W_{0,p}(s,\varphi_{\nu,p})=0$ for every $\nu_p\neq0$.
For odd $p$ this is \cite[Lemma~18(1)]{Yang} with $p\mid(DN,d)$, Gram
$\mathrm{diag}(\epsilon_1,\epsilon_2)$ with $p\,\|\,\epsilon_2$, as in the proof of
Lemma~\ref{lem:Wp}. For $p=2$, a case Yang omits, write $\lambda_0=\lambda/2=\begin{pmatrix}c_1&
c_2\\2c_3&-c_1\end{pmatrix}$ in $\mathcal O\otimes\Z_2=\begin{pmatrix}\Z_2&\Z_2\\2\Z_2&\Z_2
\end{pmatrix}$, $c_1^2+2c_2c_3=d/4\equiv2,3\bmod4$, and $L_{-,2}=\{\begin{pmatrix}a&b\\2c&-a
\end{pmatrix}:ac_1+bc_3+cc_2=0\}$ with $Q=-(a^2+2bc)$ (the reduced norm; only its integrality
is used below). If $d/4\equiv3\bmod4$ then $c_1,c_2,c_3$ are
odd; eliminating $a$ gives a form $2(uy^2+2v\,yz+u'z^2)$ with $u,u'$ units and $v$ even, hence
$2\,\mathrm{diag}(\epsilon_1,\epsilon_2)$ with $\epsilon_1\epsilon_2=-(d/4)/c_1^2\equiv1\bmod4$,
and the three nonzero cosets of $(\Z/2)^2$ have $Q\equiv\epsilon_1/4,\epsilon_2/4,\epsilon_1/2$,
none integral. If $d/4\equiv2\bmod4$ then $c_1$ is even and $c_2,c_3$ odd; eliminating $b$ gives
$2(a^2-2(c_1/c_3)ac-2(c_2/c_3)c^2)$, hence $2\,\mathrm{diag}(1,2\epsilon')$ with
$\epsilon'$ a unit, whose discriminant group $\Z/2\times\Z/4$ has $Q(a/2,b/4)=a^2/4+\epsilon'b^2/8$,
never integral on a nonzero element. This is the firing rule $\sum_{p\mid N/(N,d)}\log p$: the
primes of $N$ dividing $d$ contribute nothing, for the reason the recipe assumed at all of them.

For $m>0$ the same prime enters through the local Whittaker polynomial of the anisotropic plane,
which the routine computes from \cite[Thms.~4.3, 4.4]{KY} on the actual $L_{-,p}$; those theorems
need no unimodularity, unlike Schofer's closed form \cite[Thm.~4.1]{Schofer}, which the routine does
not use. Forcing every published point at which a level prime divides $d_0$ on the sixteen
Guo--Yang tables with $N>1$ --- $104$ points --- through the table comparison gives agreement at all
of them (\texttt{vvdata/weyl-campaign/noncoprime/noncop.m} on the campaign branch; the one
disagreement, $X_0^{21}(2)$ at $-420$, is the misprint of Remark~\ref{rem:divisor}). So the values
at such points need no separate treatment; the pipeline's coprimality filter keeps them out of the
Borcherds search only because admitting them changes its anchors and hence the forms.
\end{remark}

Two general facts used in Part~I:

\begin{lemma}[the plane's isotropy is a property of the field]\label{lem:split}
For every prime $p$, the quadratic space $\Lm\otimes\Q_p=\lambda^\perp\subset B_0\otimes\Q_p$ is
isotropic if and only if $p$ splits in $k=\Q(\sqrt{d_0})$. In particular this does not depend on
the conductor.
\end{lemma}
\begin{proof}
Write $B=k\oplus kj$ with $j$ anticommuting with $k$ and $j^2=\beta\in\Q^\times$, and $B_0$ for the
trace-zero part. Then $\lambda\in k_0=\Q\sqrt{d_0}$ and $\lambda^\perp=kj$: for $x\in k$,
$\tr(\lambda\overline{xj})=\tr(\lambda\,\bar j\,\bar x)=-\tr(\lambda\bar x j)=0$ since $\lambda\bar x
\in k$ and $\tr(yj)=0$ for $y\in k$. The quadratic form is $Q(xj)=\nrd(xj)=-\beta\,N_{k/\Q}(x)$.
So $\lambda^\perp\otimes\Q_p\cong(k_p,-\beta N)$, which represents $0$ nontrivially exactly when
$N_{k_p/\Q_p}$ does, i.e.\ when $k_p$ is not a field.
\end{proof}

\begin{lemma}[pole or no pole at the zero coset]\label{lem:pole}
In the counting normalisation of \textup{Lemma~lem:m0loc of the preprint}, $W_{0,p}(s,\varphi_0)=
(1-X)\sum_{k\ge0}\alpha_kX^k$ with $X=p^{-s}$, $\alpha_0=1$ and $\alpha_k=p^{-k}\#\{x\in
L_{-,p}/p^kL_{-,p}:\ Q(x)\equiv0\bmod p^k\}$. If $L_{-,p}\otimes\Q_p$ is anisotropic, the $\alpha_k$
are bounded and $W_{0,p}(s,\varphi_0)$ is holomorphic at $s=0$. If it is isotropic, $\alpha_k\ge
ck$ for some $c>0$ and $W_{0,p}(s,\varphi_0)$ has a pole at $s=0$.
\end{lemma}
\begin{proof}
Anisotropic: $Q$ is a nonzero anisotropic form on $\Z_p^2$, so there is $c_0$ with
$\ord_p Q(x)\le c_0+2\min_i\ord_p(x_i)$ for all $x\ne0$ (on the compact set of primitive $x$,
$\ord_pQ(x)$ is bounded, and $Q(p^ax)=p^{2a}Q(x)$). Hence $Q(x)\equiv0\bmod p^k$ forces every
coordinate of $x$ to vanish modulo $p^{\lceil(k-c_0)/2\rceil}$, which leaves at most $p^{k+c_0}$
classes modulo $p^k$, and $\alpha_k\le p^{c_0}$. Then $\lvert(1-X)\sum\alpha_kX^k\rvert\le p^{c_0}$
for $0<X<1$, so there is no pole. Isotropic: choose a primitive isotropic $e\in L_{-,p}$ and
complete it to a basis $e,f'$, so $Q(xe+yf')=uxy+wy^2$ with $u=\langle e,f'\rangle\ne0$. For each
$0\le j\le k-\ord_pu-1$, the classes $y$ with $\ord_py=j$ number $(p-1)p^{k-j-1}$, and for each of
them $Q\equiv0\bmod p^k$ is a congruence on $x$ modulo $p^{k-j-\ord_pu}$ with $p^{j+\ord_pu}$
solutions modulo $p^k$; summing, $\alpha_k\ge(k-\ord_pu)(p-1)/p$, which grows linearly, and the
pole follows since $(1-X)\sum\alpha_kX^k\ge(1-X)\sum_{k\le K}ckX^k$ for every $K$.
\end{proof}

\section{Non-fundamental discriminant}
Write $d=\ell^2d_0$ as in Part~I (the standalone writes $f$ for the conductor; here $f$ is the input form), $R_\ell=\Z+\ell\mathcal O_k$, and let $\lambda$ be the image of $\sqrt d$ under an optimal embedding of $R_\ell$.

\begin{lemma}[the lattice at a conductor prime]\label{lem:conductor}
Let $p$ be an odd level prime, $p\,\|\,N$, $p\nmid d_0$, $p^k\,\|\,\ell$ with $k\ge1$, and
$\lambda$ the image of $\sqrt d=f\sqrt{d_0}$ under an optimal embedding of $R_\ell$. In the model
$\mathcal O\otimes\Z_p=\begin{pmatrix}\Z_p&\Z_p\\p\Z_p&\Z_p\end{pmatrix}$, $L_p=\{\begin{pmatrix}
a&b\\pc&-a\end{pmatrix}\}$ with $Q=\nrd$, so $Q(a,b,c)=-a^2-pbc$:
\begin{enumerate}
\item $\lambda=\begin{pmatrix}pa'&b\\p^2c'&-pa'\end{pmatrix}$ with $b$ a unit and
$a'^2+bc'=-\lvert d\rvert/p^2$;
\item $L_{-,p}$ has the basis $e_1=(1,0,-2a'/b)$, $e_2=(0,1,-pc'/b)$ and, after the integral change
of basis $e_1'=e_1$, $e_2'=e_2-(pa'/b)e_1$, the Gram form
$Q=-\alpha^2-p^{2k}\lvert d_0\rvert v^2\beta^2$ with $v$ a unit: a unit line orthogonal to a
$p^{2k}$-scaled line;
\item $[L_p:\Z_p\lambda\oplus L_{-,p}]=p^{2k-1}$;
\item $L_{-,p}^\vee/L_{-,p}\cong\Z/p^{2k}$, generated by $e_2'/p^{2k}$; its cosets of integral norm are
the $p^k-1$ multiples of $e_2'/p^k$, and of these exactly the $p-1$ multiples of $e_2'/p$ lie in
$L_p^\vee$. In particular the direct-sum decomposition of Lemma~\ref{lem:LN} fails, and a coset of
$L_{-,p}^\vee/L_{-,p}$ enters the decomposition of some $\varphi_\eta$, $\eta\in\Ldual/L$, only if it
is one of those $p-1$; the others can pair only with $x\notin\Ldual$, $x+\nu\in\Ldual$
(Theorem~\ref{thm:kappa}(c)).
\end{enumerate}
\end{lemma}
\begin{proof}
Write $\lambda=(a,b,c)$. Then $Q(\lambda)=-a^2-pbc=\lvert d\rvert\cdot(\text{unit})^2$ is divisible by
$p^{2k}$ with $2k\ge2$, so $p\mid a$. Optimality of the embedding of $R_\ell$ says $\lambda/p\notin
\mathcal O\otimes\Z_p$, i.e.\ not all of $a,b,c$ are divisible by $p$, so $b$ or $c$ is a unit.
Conjugation by the Atkin--Lehner element $w_p=\begin{pmatrix}0&1\\p&0\end{pmatrix}$ normalises
$\mathcal O\otimes\Z_p$, preserves the norm, and sends $(a,b,c)$ to $(-a,c,b)$; since everything
below is an isometry invariant of the pair $(L_p,\lambda)$, we may take $b$ to be the unit. Then
$p^2\mid a^2$ forces $p\mid c$, which is (1). Conjugating further by $\begin{pmatrix}1&0\\pt&1
\end{pmatrix}\in(\mathcal O\otimes\Z_p)^\times$ replaces $a'$ by $a'-tb$, so we may and do take
$a'=0$ below; this is what makes the pairing $\langle e_2',x\rangle=-p\gamma+(p^2c'/b)\beta$ free of an
$\alpha$-term in the proof of (4). For $x=(\alpha,\beta,\gamma)\in L_p$,
$\langle x,\lambda\rangle=-\tr(x\lambda)=-(2\alpha pa'+\beta p^2c'+\gamma pb)$, so $x\perp\lambda$
iff $\gamma=-(2a'\alpha+pc'\beta)/b$, which gives the basis in (2); substituting,
$Q|_{L_{-,p}}=-\alpha^2+(2pa'/b)\alpha\beta+(p^2c'/b)\beta^2$, and completing the square gives
$-\alpha'^2+(p^2/b^2)(a'^2+bc')\beta^2=-\alpha'^2-p^{2k}\lvert d_0\rvert v^2\beta^2$ by (1). For
(3) the determinant of the matrix with rows $\lambda,e_1,e_2$ is $(2p/b)(a'^2+bc')$, of valuation
$2k-1$. For (4), $\langle e_2',e_2'\rangle=-2p^{2k}\lvert d_0\rvert v^2$ and $\langle e_1',e_1'\rangle
=-2$ are the only nonzero pairings, so the dual is spanned by $e_1'/2$ and $e_2'/(2p^{2k}\lvert d_0
\rvert v^2)$, i.e.\ by $e_2'/p^{2k}$ up to units; $Q(je_2'/p^{2k})=-\lvert d_0\rvert v^2j^2/p^{2k}$ is
integral iff $p^k\mid j$. Finally $\langle e_2'/p^m,x\rangle=-(1/p^{m-1})(\gamma+(pc'/b)\beta)+
\dots$ is integral on $L_p$ iff $m\le1$, which leaves the multiples of $e_2'/p$.
\end{proof}

\begin{lemma}[the lattice at the conductor prime $2$]\label{lem:conductor2}
Let $2\,\|\,N$, $2\nmid d_0$, $2^k\,\|\,\ell$ with $k\ge1$, and $\lambda$ the image of $\sqrt d$ under an
optimal embedding of $R_\ell$, in the model $\mathcal O\otimes\Z_2=\begin{pmatrix}\Z_2&\Z_2\\2\Z_2&\Z_2
\end{pmatrix}$ with $Q=\nrd$, $Q(a,b,c)=-a^2-2bc$ for $\begin{pmatrix}a&b\\2c&-a\end{pmatrix}$.
\begin{enumerate}
\item After conjugating by $w_2$ if necessary, $\lambda=(2a',2b',2c')$ with $b'$ a unit, $c'$ even
      and $a'^2+2b'c'=-4^{k-1}\lvert d_0\rvert$ ($a'$ odd exactly when $k=1$); $\lambda_0=\lambda/2$
      is primitive in $L_2$ and $L_{+,2}=\Z_2\lambda_0$.
\item $L_{-,2}$ has the basis $e_1=(1,0,-a'/b')$, $e_2'=(0,1,-c'/b')-(a'/b')e_1$ with
      $Q=-\alpha^2-4^{k-1}c\,\beta^2$, $c=\lvert d_0\rvert/b'^2\equiv3\bmod4$, and $-c\equiv1\bmod8$
      exactly when $2$ splits in $k$.
\item $[L_2:L_{+,2}\oplus L_{-,2}]=2^{2k-2}$; in particular the sum is direct at conductor $2$.
\item $L_{-,2}^\vee/L_{-,2}\cong\Z/2\times\Z/2^{2k-1}$, generated by $e_1/2$ and $e_2'/2^{2k-1}$. Its
      $2^k$ cosets of integral norm are $(0,t/2^{k-1})$, $t\bmod2^{k-1}$, and $(1/2,t/2^k)$, $t$ odd
      mod $2^k$; writing $2^{k-\rho}$ for the order of a nonzero one, $\rho=1+\ord_2t\in[1,k-1]$ for
      the first kind and $\rho=0$ for the second, so there are $2^{k-\rho-1}$ cosets of each
      $\rho\in[0,k-1]$. A nonzero coset lies in $L_2^\vee$ iff it has order $2$ ($\rho=k-1$); there is
      exactly one.
\end{enumerate}
\end{lemma}
\begin{proof}
Over $\Z_2$, $R_\ell\otimes\Z_2=\Z_2+2^k\mathcal O_{k}$ is generated by $2^{k-1}(1+\sqrt{d_0})$, whose
image is $2^{k-1}+\lambda/(2\ell')$ with $\ell=2^k\ell'$, $\ell'$ a unit; it lies in $\mathcal O\otimes\Z_2$ iff
$\lambda/2$ does, i.e.\ iff $a,b,c$ are all even. Optimality says the order of conductor $2^{k-1}$
does not embed: for $k\ge2$ its generator $2^{k-2}(1+\sqrt{d_0})$ maps to $2^{k-2}+\lambda/(4\ell')$, so
$\lambda/4\notin\mathcal O\otimes\Z_2$, i.e.\ not all of $a',b',c'$ are even; for $k=1$ the generator
$(1+\sqrt{d_0})/2$ maps to $(1+\lambda/(2\ell'))/2$, which lies in the order iff $a'$ is odd and $b',c'$
are even. From $Q(\lambda)=4^k\lvert d_0\rvert$ we get $a'^2+2b'c'=-4^{k-1}\lvert d_0\rvert$. If
$k\ge2$ the right side is $\equiv0\bmod4$, which forces $a'$ even and then $b'c'$ even, so by
optimality exactly one of $b',c'$ is odd; if $k=1$ the right side is odd, so $a'$ is odd,
$2b'c'=-\lvert d_0\rvert-a'^2\equiv0\bmod4$ makes $b'c'$ even, and optimality again says one of
$b',c'$ is odd. Conjugation by $w_2$ sends $(a',b',c')$ to $(-a',c',b')$, so we may take $b'$ odd and
$c'$ even. Then $\lambda_0=(a',b',2c')$ has an odd entry, hence is primitive. This is (1).

For (2): $\langle x,\lambda\rangle=-2(a\alpha+c\beta+b\gamma)=-4(a'\alpha+c'\beta+b'\gamma)$ for
$x=(\alpha,\beta,\gamma)$, so $x\perp\lambda$ iff $\gamma=-(a'\alpha+c'\beta)/b'$, which gives the
basis $e_1,e_2=(0,1,-c'/b')$ with $Q(\alpha e_1+\beta e_2)=-\alpha^2+(2a'/b')\alpha\beta+(2c'/b')
\beta^2$; completing the square (an integral change of basis since $b'$ is a unit) gives $-\alpha'^2
+((a'^2+2b'c')/b'^2)\beta^2=-\alpha'^2-4^{k-1}c\,\beta^2$. As $b'$ is odd, $c\equiv\lvert d_0\rvert
\bmod8$, and $d_0\equiv1\bmod4$ gives $\lvert d_0\rvert\equiv3\bmod4$, with $\lvert d_0\rvert\equiv7
\bmod8$ iff $d_0\equiv1\bmod8$ iff $2$ splits.

For (3): the Gram matrix of $\langle\cdot,\cdot\rangle=2Q$ on $L_2$ in the coordinates $(\alpha,
\beta,\gamma)$ is $(-2)\oplus\begin{pmatrix}0&-2\\-2&0\end{pmatrix}$, of determinant $8$; on
$L_{+,2}\oplus L_{-,2}$ it is $\mathrm{diag}(2\cdot4^{k-1}\lvert d_0\rvert,\,-2,\,-2\cdot4^{k-1}c)$,
of determinant $8\cdot4^{2k-2}\cdot$unit. The ratio of determinants is the square of the index.

For (4): the dual of $\mathrm{diag}(-2,-2^{2k-1}c)$ is spanned by $e_1/2$ and $e_2'/(2^{2k-1}c)$, and
$Q(se_1/2+te_2'/2^{2k-1})=-s^2/4-ct^2/2^{2k}$. For $s=0$ this is integral iff $2^k\mid t$; for $s=1$
iff $ct^2\equiv-2^{2k-2}\bmod2^{2k}$, i.e.\ $t=2^{k-1}t_1$ with $t_1$ odd and $ct_1^2\equiv-1\bmod4$,
which holds since $c\equiv3\bmod4$. The orders are read off the denominators. Finally
$\langle\nu,x\rangle=-2(\nu_\alpha\alpha+\nu_\beta\gamma+\nu_\gamma\beta)$, so $\nu\in L_2^\vee$ iff
$2\nu_\alpha,2\nu_\beta,2\nu_\gamma\in\Z_2$. For $\nu=ue_1+we_2'$ one has $\nu_\beta=w$, so
$2w\in\Z_2$ is necessary: for the first kind $w=t/2^{k-1}$ this forces $2^{k-2}\mid t$, i.e.\
$\rho=k-1$, and for the second kind $w=t/2^k$ with $t$ odd it forces $k=1$, again $\rho=k-1$; in both
cases $2\nu_\alpha=2u-2(a'/b')w$ and $2\nu_\gamma=-2(a'/b')u+2w(a'^2/b'^2-c'/b')$ are then integral.
\end{proof}

\begin{lemma}[the $m=0$ local factors at a conductor prime]\label{lem:Wcond}
Let $p\,\|\,N$, $p\nmid d_0$, $p^k\,\|\,\ell$ with $k\ge1$, and $\epsilon=(\tfrac{d_0}{p})=\pm1$,
with $L_{-,p}$ as in Lemma~\ref{lem:conductor} ($p$ odd) or Lemma~\ref{lem:conductor2} ($p=2$). In
the counting normalisation $W_{0,p}(s,\varphi_\nu)=(1-X)\sum_{j\ge0}
\alpha_j(\nu)X^j$, $X=p^{-s}$, $\alpha_j(\nu)=p^{-j}\#\{x\in(\nu+L_{-,p})/p^jL_{-,p}:Q(x)\equiv0
\bmod p^j\}$:
\begin{enumerate}
\item For $p$ odd the cosets of integral norm are $\nu_r=(0,r/p^k)$, $r\bmod p^k$; for
      $r\neq0$ let $\rho=\ord_p r<k$, so that $\nu_r$ has order $p^{k-\rho}$ and lies in
      $L_p^\vee$ iff $\rho=k-1$. For $p=2$ let $\nu_r$ run over the $2^k-1$ nonzero integral cosets
      of Lemma~\ref{lem:conductor2}(4), with $\rho$ defined there by the order $2^{k-\rho}$. In
      both cases $\alpha_j(\nu_r)=p^{\lfloor j/2\rfloor}$ for $j\le2\rho$ and $\alpha_j(\nu_r)=
      (1+\epsilon)p^\rho$ for $j>2\rho$, i.e.
      \[
        W_{0,p}(s,\varphi_{\nu_r})=(1-X)\sum_{j=0}^{2\rho}p^{\lfloor j/2\rfloor}X^j
                                   +(1+\epsilon)\,p^\rho X^{2\rho+1}.
      \]
\item For the zero coset, $\alpha_j(0)=p^{\lfloor j/2\rfloor}$ for $j\le2k$ and
      $\alpha_j(0)=(1+\epsilon)(p-1)p^{k-1}\lceil(j-2k)/2\rceil+p^{\,k-(j\bmod2)}$ for $j>2k$, i.e.
      \[
        W_{0,p}(s,\varphi_0)=(1-X)\Bigl[\sum_{j=0}^{2k}p^{\lfloor j/2\rfloor}X^j
          +\frac{X^{2k+1}(p^{k-1}+p^kX)}{1-X^2}\Bigr]
          +\frac{(1+\epsilon)(p-1)p^{k-1}X^{2k+1}}{1-X^2}.
      \]
\item Consequently $\kappa^-_{\nu_r}(0)=-\dfrac{2p^{\rho-k+1}}{p-1}\log p$ if $p$ splits, and
      \[
        \kappa^-_{\nu_r}(0)=-\frac{2}{p^{k-1}(p+1)}\Bigl(p^\rho+\frac{2(p^\rho-1)}{p-1}\Bigr)\log p
      \]
      if $p$ is inert; all other nonzero cosets have $\kappa^-_\nu(0)=0$.
\end{enumerate}
\end{lemma}
\begin{proof}
With $Q(u,t)=-(u^2+c\,p^{2k}t^2)$ the dual is generated by $e_2/p^{2k}$, and $Q(0,r/p^{2k})=
-cr^2/p^{2k}$ is integral iff $p^k\mid r$, which is (1)'s list and Lemma~\ref{lem:conductor}(4).
A vector of $\nu_r+L_{-,p}$ is $(u,w)$ with $w=r+p^kt\in r+p^k\Z_p$, and for $r\ne0$ every such $w$
has $\ord_pw=\rho$. The congruence $u^2+cw^2\equiv0\bmod p^j$ reads $u^2\equiv0$ when $j\le2\rho$,
with $p^{\lfloor j/2\rfloor}$ solutions $u\bmod p^j$ for each of the $p^j$ values of $t$; when
$j>2\rho$ it forces $u=p^\rho u'$ with $u'^2\equiv-c(w/p^\rho)^2\bmod p^{j-2\rho}$, which has $2$
solutions $u'\bmod p^{j-2\rho}$ if $-c$ is a square mod $p$ --- i.e.\ $\epsilon=1$, since
$-c=d_0v^2$ --- and none otherwise, giving $(1+\epsilon)p^{\rho}$ values of $u\bmod p^j$ for each
$t$. This is (1). For the zero coset, $w=p^kt$; if $t\equiv0\bmod p^j$ or $2k+2\ord_pt\ge j$ the
congruence is $u^2\equiv0$ with $p^{\lfloor j/2\rfloor}$ solutions, and the $t$ with
$2k+2v<j$, $v=\ord_pt$, number $(p-1)p^{j-v-1}$ and each admits $(1+\epsilon)p^{k+v}$ values of
$u$ by the same Hensel count. Summing, for $j>2k$,
\[
  p^j\alpha_j(0)=(1+\epsilon)(p-1)p^{j+k-1}\Bigl\lceil\frac{j-2k}2\Bigr\rceil
   +p^{\lfloor j/2\rfloor}\Bigl(1+\sum_{v\ge\lceil(j-2k)/2\rceil}^{j-1}(p-1)p^{j-1-v}\Bigr)
   =(1+\epsilon)(p-1)p^{j+k-1}\Bigl\lceil\frac{j-2k}2\Bigr\rceil+p^{\lfloor j/2\rfloor+j-\lceil(j-2k)/2\rceil},
\]
and $\lfloor j/2\rfloor-\lceil (j-2k)/2\rceil=k-(j\bmod2)$; for $j\le2k$ only $u^2\equiv0$ remains.
The generating functions follow from $\sum_{n\ge1}\lceil n/2\rceil X^n=X/((1-X)(1-X^2))$. For (3)
apply Theorem~\ref{thm:kappa}(c). If $\epsilon=-1$, $W_{0,p}(s,\varphi_0)$ is finite at $X=1$ with
value $(p^{k-1}+p^k)/2$ (only the middle term survives the factor $1-X$) and $W_{0,p}(s,
\varphi_{\nu_r})=(1-X)\sum_{j\le2\rho}p^{\lfloor j/2\rfloor}X^j$ has derivative
$-\sum_{j\le2\rho}p^{\lfloor j/2\rfloor}=-(p^\rho+2(p^\rho-1)/(p-1))$ there. If $\epsilon=1$,
$(1-X)W_{0,p}(s,\varphi_0)\to(p-1)p^{k-1}$ at $X=1$ while $W_{0,p}(s,\varphi_{\nu_r})\to2p^\rho$,
so the ratio has derivative $-2p^\rho/((p-1)p^{k-1})$.

For $p=2$ the form is $u^2+4^{k-1}cw^2$ with $c\equiv3\bmod4$ and $-c\equiv1\bmod8$ iff
$\epsilon=1$ (Lemma~\ref{lem:conductor2}), and the count of $u'$ modulo $2^n$ with $u'^2\equiv
-c\,w'^2$, $w'$ odd, is $1$, $2$, $2(1+\epsilon)$ for $n=1$, $n=2$, $n\ge3$. For the zero coset
and for the cosets $(0,t/2^{k-1})$ write the condition as $u^2+cW^2\equiv0\bmod2^j$ with
$W=2^{k-1}w$, so that $\ord_2W=k-1+\ord_2w$ for the zero coset and $\ord_2W=\rho-1$ for the coset
$(0,t/2^{k-1})$ of depth $\rho$. With $v=\ord_2W$ and $n=j-2v$: for $n\le0$ the condition is
$u^2\equiv0$, with $2^{\lfloor j/2\rfloor}$ solutions; for $n=1,2$ it is $u=2^vu'$ with $1$,
resp.\ $2$ values of $u'$ modulo $2^n$, i.e.\ $2^{v}$, resp.\ $2^{v+1}$ values of $u$ modulo $2^j$,
which again equals $2^{\lfloor j/2\rfloor}$ since $j=2v+1,2v+2$; for $n\ge3$ it is $(1+\epsilon)
2^{v+1}$ values. These are exactly the per-$w$ counts of the odd case with $p^{k+v}$ replaced by
$2^{(k-1+\ord_2 w)+1}$, so the sums are the same and (1), (2) hold verbatim: for the coset of depth
$\rho\ge1$ the count is $2^{\lfloor j/2\rfloor}$ for $j\le2\rho$ and $(1+\epsilon)2^\rho$ beyond.
For a coset $(1/2,t/2^k)$, $t$ odd, put $U=2u$ and $W=2^kw$, both odd; the condition becomes
$U^2+cW^2\equiv0\bmod2^{j+2}$, which holds for every $U$ when $j=0$ (as $1+c\equiv0\bmod4$) and for
$j\ge1$ has $2(1+\epsilon)$ solutions $U$ modulo $2^{j+2}$, i.e.\ $(1+\epsilon)$ solutions $u$
modulo $2^j$ for each of the $2^j$ values of $w$: $\alpha_j=1+\epsilon$ for $j\ge1$, the case
$\rho=0$. Part (3) follows as before.
\end{proof}

The formulas were also checked by brute-force counting on the plane $L\cap\lambda^\perp$ of
$X_0^{10}(3)$ at $p=3$ --- $d=-72,-648$ (split, $k=1,2$) and $d=-387,-603$ (inert, $k=1$), every
$3$-primary integral coset, six terms each (\texttt{level-p2/wcond\_check.m}). At $p=2$ the
six rows of the table are the lemma's series for $k=1,2$ and both splitting types, and the factors
$-2$, $-1$, $-\tfrac43$, $-\tfrac13$, $-\tfrac23\log2$ are the values of (3) at
$(k,\rho)=(1,0),(2,1),(2,0)$ in the split case and $(2,1),(2,0),(1,0)$ in the inert case. With
Lemma~\ref{lem:Wcond}, Theorem~\ref{thm:kappa}(c) is explicit at a conductor prime: the restored
terms there are
\[
  -\frac14\sum_{r\neq0}\kappa^-_{\nu_r}(0)\sum_{\substack{x\in\Lp^\vee\\ x+\nu_r\in\Ldual}}
  c_{[x+\nu_r]}\bigl(-Q(x)\bigr),
\]
with the fibre over each $\nu_r$ read off the lattice as in the computation above. The routine
evaluates the sum as written at a level prime dividing the conductor (\texttt{M0FibreCorrection}),
with the closed-form $\kappa^-_\nu(0)$ and the fibre read off the lattice, and reproduces the $18$
values of the table together with the $28$ at conductors $1$ and $2$ (\texttt{tests/M0PoleSum.m}).

\subsection*{A prime of the conductor outside the level}

A prime $q$ dividing the conductor need not divide the level ($q\mid D$ is impossible: at a prime
of $D$ the only order of $k_q$ embedding optimally into the maximal order of the division algebra is
the maximal one). There the same thing happens with one difference, and it puts terms in $\log q$
into the formula at a prime where the lattice $L$ itself is unimodular.

\begin{lemma}[the lattice at an odd conductor prime outside the level]\label{lem:unimod}
Let $q$ be odd, $q\nmid DN$, $q^k\,\|\,\ell$ with $k\ge1$, and $\lambda$ the image of $\sqrt d$ under
an optimal embedding of $R_\ell$. In $\mathcal O\otimes\Z_q=M_2(\Z_q)$, after conjugating by
$GL_2(\Z_q)$, $\lambda=\begin{pmatrix}0&d\\1&0\end{pmatrix}$, so $L_{+,q}=\Z_q\lambda$ and
\begin{enumerate}
\item $L_{-,q}=\langle-1\rangle\perp\langle-\lvert d\rvert\rangle$, with basis $e_1=\mathrm{diag}
      (1,-1)$, $e_2=\begin{pmatrix}0&-d\\1&0\end{pmatrix}$; $[L_q:L_{+,q}\oplus L_{-,q}]=q^{2k}$;
\item $L_{-,q}^\vee/L_{-,q}$ has $q$-part $\Z/q^{2k}$ (if $q\nmid d_0$) or $\Z/q^{2k+1}$ (if
      $q\mid d_0$), generated by $e_2/\lvert d\rvert$, and in both cases its cosets of integral norm are
      the $q^k$ cosets $\nu_r=(0,r/q^k)$, $r\bmod q^k$, of order $q^{k-\rho}$ with $\rho=\ord_qr$;
      \emph{none} of the nonzero ones lies in $L_q^\vee=L_q$;
\item the fibre over $\nu_r$ is $\{t\lambda:\ t\in-r/q^k+\Z_q\}\cap\Lp^\vee$: the vectors
      $x=t\lambda$ with $t\equiv-r/q^k\bmod\Z_q$, whose norms $t^2\lvert d\rvert$ are those of the
      primitive CM vectors of the discriminants $d/q^{2(k-\rho)}$ and their multiples prime to $q$;
      the class $[x+\nu_r]\in\Ldual/L$ is trivial at $q$ and equals that of $x$ elsewhere, so for
      $x\in L[1/q]$ it is $0$ and the coefficient is $c_0(-Q(x))$;
\item if $q\nmid d_0$ the plane is $\langle-1\rangle\perp\langle-q^{2k}c\rangle$ with $c$ a unit and
      Lemma~\ref{lem:Wcond} applies verbatim; if $q\mid d_0$ it is $\langle-1\rangle\perp\langle
      -q^{2k+1}c'\rangle$, anisotropic, with $\alpha_j(\nu_r)=q^{\lfloor j/2\rfloor}$ for
      $j\le2\rho+1$ and $0$ beyond, $\alpha_j(0)=q^{\lfloor j/2\rfloor}$ for $j\le2k+1$ and $q^k$
      beyond, hence
      \[
        \kappa^-_{\nu_r}(0)=-\frac{2(q^{\rho+1}-1)}{(q-1)q^k}\log q .
      \]
\end{enumerate}
Consequently the restored terms at $q$ are, per CM point,
\[
  -\frac14\sum_{\rho=0}^{k-1}\kappa^-_\rho(0)\sum_{\substack{x\in q^{\rho-k}\Lambda\setminus q^{\rho-k+1}\Lambda}}
   c_{[x]'}\bigl(-Q(x)\bigr),\qquad \Lambda=\Ldual\cap\Q\lambda,
\]
with $[x]'$ the class of $x$ away from $q$; when every such $x$ lies in $L[1/q]$ this is
$-\tfrac12\sum_{j=1}^{k}\kappa^-_{k-j}(0)\sum_{m\ge1,\ q\nmid m}c_0\bigl(-m^2Q(\lambda_0)/q^{2j}\bigr)$,
a pole sum over the CM points of the lower discriminants $d/q^2,\dots,d/q^{2k}$.
\end{lemma}
\begin{proof}
Since $q^{2k}\mid Q(\lambda)=\lvert d\rvert$ and $\lambda$ is primitive (optimality: $\lambda/q\notin
\mathcal O_q$, as $R_\ell\otimes\Z_q=\Z_q+q^k\mathcal O_{k,q}$ is generated by $q^k\sqrt{d_0}=
\lambda/\ell'$ with $\ell'$ a unit), $\lambda$ reduces mod $q$ to a nonzero nilpotent matrix, of rank
$1$; a vector $v$ with $\lambda v\not\equiv0$ then gives the basis $v,\lambda v$ of $\Z_q^2$, in
which $\lambda=\begin{pmatrix}0&d\\1&0\end{pmatrix}$. With $Q(a,b,c)=-a^2-bc$ on $\begin{pmatrix}
a&b\\c&-a\end{pmatrix}$, $\langle x,\lambda\rangle=-(\beta+d\gamma)$, which gives the basis of
$L_{-,q}$ and its Gram matrix; the determinants $2$ of $L_q$ and $-8\lvert d\rvert^2$ of $L_{+,q}
\oplus L_{-,q}$ give the index. The dual is generated by $e_2/\lvert d\rvert$, and with $\lvert d\rvert=q^{2k+\delta}c$
($\delta\in\{0,1\}$, $c$ a unit) the $q$-part is generated by $e_2/q^{2k+\delta}$; $Q(te_2/q^{2k+\delta})
=-ct^2/q^{2k+\delta}$ is integral iff $q^{k+\delta}\mid t$, i.e.\ iff the coset is a multiple of
$e_2/q^k$. No nonzero $\nu_r$ lies in $L_q$, which is $L_q^\vee$.
For (3), $x=t\lambda=(0,td,t)$ and $\nu_r=(0,-dr/q^k,r/q^k)$, so $x+\nu_r\in L_q$ iff $t+r/q^k
\in\Z_q$ (the other coordinate then differs from $d(t+r/q^k)$ by $2dr/q^k\in\Z_q$); at $p\neq q$
the condition is $x\in L_p^\vee$, and $\nu_r$ has no component there. For (4), the odd case is
Lemma~\ref{lem:Wcond}; in the ramified case write $W=q^kw$, so $Q=-(u^2+qc'W^2)$ with $\ord_q(qc'W^2)$
odd, and $u^2\equiv-qc'W^2\bmod q^j$ has solutions only when both sides are $\equiv0$, i.e.\ for
$j\le\ord_q(qc'W^2)$, with $q^{\lfloor j/2\rfloor}$ of them; summing over $W$ as in the proof of
Lemma~\ref{lem:Wcond} gives the series, which are bounded (no pole, as the plane is anisotropic), and
Theorem~\ref{thm:kappa}(c) gives $\kappa^-=-\sum_{j\le2\rho+1}q^{\lfloor j/2\rfloor}/q^k\cdot\log q$.
The display is Theorem~\ref{thm:kappa}(c) with the fibres of (3), the cosets of depth $\rho$
covering the residues $t\equiv-r/q^k$ with $\ord_qr=\rho$ exactly once.
\end{proof}

\paragraph{Check.} Table~45 does not reach such points, but the forms must still be functions of
the one Hauptmodul $s$, so at $d=-588=14^2\cdot(-3)$ on $X_0^{15}(2)$ ($7$ splits in $\Q(\sqrt{-3})$,
three star points) the nine values (now sums over the three points) determine $N(s)$, $N(s-2)$ and
$N((s+\tfrac1{12})(s-\tfrac54))$, and a monic cubic $H\in\Q[X]$ with these absolute values at $0$,
$2$ and (as a product) at $-\tfrac1{12},\tfrac54$ must exist --- a rationality condition on the
discriminant of a quadratic. With the term of Lemma~\ref{lem:unimod} (per point
$\tfrac16c_\infty(-3)\log7$, the pole of $f$ at $\tau_{-12}$) it exists,
\[
  H(X)=X^3-\tfrac{191}{54}X^2+\tfrac{343}{432}X-\tfrac{83^2}{2^8\cdot3^3\cdot7},
\]
and the cubic field it cuts out has discriminant $-588$ and the same reduced polynomial
$X^3-X^2+5X+1$ as the field of definition of the CM points computed by class field theory
(\texttt{FieldsOfDefinitionOfCMPoint}); without the term no rational cubic exists. The $7$-adic
Newton polygon of $H$ puts every root at $\ord_7=-\tfrac13$: all three points reduce to the pole
$\tau_{-12}$, which is the congruence $s(\tau_{dq^2})\equiv s(\tau_d)$ at the primes above $q$ that
the cyclic $q$-isogeny between the CM points predicts. The same happens at $d=-1960=7^2\cdot(-40)$
(field discriminant $-1960$, roots divisible by $7$ as $s(\tau_{-40})=0$), and at $d=-735=7^2\cdot
(-15)$, $7$ inert, the term puts the factor $7$ into $N(s-\tfrac54)$ as $s(\tau_{-15})=\tfrac54$
requires (\texttt{level-p2/qprime.m}, \texttt{classpoly.py}, \texttt{fod.m}). Three published values do
reach such points, once they are forced into the table (the pipeline's candidate search offers
conductors $1$ and $2$ only): on $X_0^{58}(1)$ and $X_0^{55}(1)$ --- no level prime, so no other
$m=0$ term at all --- \cite{GY} give $s(-27)=3/19$ and $0$, with $-27=3^2\cdot(-3)$ and $3\mid
d_0$, the ramified case of (4); on $X_0^{21}(2)$ they give $s(-100)=1/5$, with $-100=5^2\cdot(-4)$
and $5$ split in $\Q(i)$, the case of Lemma~\ref{lem:Wcond}. The terms are $\log3$, $-\log3$ and
$-4\log5$ on the Hauptmodul form, and with them the pipeline reproduces all three published values
(by cross-ratio against three other points of each table, as in
\texttt{tests/\_offline/GuoYangCheck.m}; \texttt{level-p2/gyforce\_*.m}). 

\begin{lemma}[the conductor prime $2$ outside the level]\label{lem:unimod2}
Let $2\nmid DN$, $2^k\,\|\,\ell$ with $k\ge1$, and $\lambda$ the image of $\sqrt d$ under an optimal
embedding of $R_\ell$. In $\mathcal O\otimes\Z_2=M_2(\Z_2)$, after conjugating by $GL_2(\Z_2)$,
$\lambda_0=\lambda/2=\begin{pmatrix}0&d/4\\1&0\end{pmatrix}$ is primitive, $L_{+,2}=\Z_2\lambda_0$,
and with $e_1=\mathrm{diag}(1,-1)$, $e_2=\begin{pmatrix}0&-d/4\\1&0\end{pmatrix}$:
\begin{enumerate}
\item $L_{-,2}=\langle-1\rangle\perp\langle-\lvert d\rvert/4\rangle$. If $d_0$ is odd this is
      the plane $\langle-1\rangle\perp\langle-4^{k-1}c\rangle$, $c\equiv3\bmod4$, of
      Lemma~\ref{lem:conductor2}, with the same $2^k$ integral cosets $(0,t/2^{k-1})$ and
      $(1/2,t/2^k)$ and the same series and $\kappa^-_\nu(0)$ (Lemma~\ref{lem:Wcond}); if $d_0$ is
      even it is $\langle-1\rangle\perp\langle-4^kc\rangle$ with $c\equiv1\bmod4$ or $\langle-1
      \rangle\perp\langle-2^{2k+1}c\rangle$, anisotropic, with integral cosets $(0,r/2^k)$,
      $\rho=\ord_2r$, series $\alpha_j=2^{\lfloor j/2\rfloor}$ for $j\le2\rho+1$ and $0$ beyond
      (zero coset: up to $j\le2k+1$, then $2^k$), and
      $\kappa^-_{\nu_r}(0)=-2(2^{\rho+1}-1)/2^k\cdot\log2$, the formula of Lemma~\ref{lem:unimod}(4)
      at $q=2$.
\item $L_2^\vee=L_2+\Z_2e_1/2$ meets $L_{-,2}\otimes\Q$ in $L_{-,2}+\Z_2e_1/2$, and $e_1/2$ has
      norm $-\tfrac14$: no nonzero integral coset of the plane lies in $L_2^\vee$. The cosets of the
      second kind $(1/2,t/2^k)$ pair with $x$ such that $x+\nu\in e_1/2+L_2$, i.e.\ over the
      nonzero class of $(\Ldual/L)_2\cong\Z/2$, at the exponents $Q(x)\equiv\tfrac34\bmod\Z$; all
      others pair over the trivial class at $2$.
\end{enumerate}
\end{lemma}
\begin{proof}
$R_\ell\otimes\Z_2=\Z_2+2^k\mathcal O_{k,2}=\Z_2[2^{k-1}\sqrt{d_0}]=\Z_2[\lambda/2]$ whether $d_0$ is
odd or even, so the embedding is optimal iff $\Z_2^2$ is a module over $\Z_2[\lambda/2]$ with that
order as its full ring of multipliers; a lattice in a quadratic field whose multiplier ring is a
given order is an invertible, hence principal, ideal of it, so $\Z_2^2\cong\Z_2[\lambda/2]$ and in
the basis $1,\lambda/2$ the element $\lambda/2$ acts as $\begin{pmatrix}0&(\lambda/2)^2\\1&0
\end{pmatrix}$ with $(\lambda/2)^2=d/4$. With $Q(a,b,c)=-a^2-bc$, $\langle x,\lambda_0\rangle=
-(\beta+\gamma d/4)$ gives the basis $e_1,e_2$ of $L_{-,2}$ and $Q(e_1)=-1$, $Q(e_2)=d/4$,
$\langle e_1,e_2\rangle=0$. For odd $d_0$, $\lvert d\rvert/4=4^{k-1}\lvert d_0\rvert \ell'^2$ with
$\lvert d_0\rvert\equiv3\bmod4$; for even $d_0=4m$, $\lvert d\rvert/4=4^k\lvert m\rvert \ell'^2$ with
$\lvert m\rvert\equiv1\bmod4$ or even, and $-c$ is then not a square mod $4$, resp.\ the exponent
is odd, so the counts of Lemma~\ref{lem:Wcond} stop after $n=1$ exactly as in the ramified case
of Lemma~\ref{lem:unimod}. The Gram matrix of $\langle\cdot,\cdot\rangle$ on $L_2$ in the basis
$e_1,e_{12},e_{21}$ is $(-2)\oplus\begin{pmatrix}0&-1\\-1&0\end{pmatrix}$, so $L_2^\vee=L_2+\Z_2
e_1/2$; a vector $\nu=se_1/2+te_2/(\lvert d\rvert/2)$ has $\langle\nu,e_{21}\rangle=-t/2$ and
$\langle\nu,e_{12}\rangle=-2t/\lvert d\rvert$, integral only for $t\equiv0$. For (2), $x=t\lambda_0
=(0,td/4,t)$ and $x+\nu\in L_2^\vee$ iff $t+2t'/\lvert d\rvert\in\Z_2$ (then the $\beta$-coordinate
is integral automatically), for either value of $s$; when $s=1$ the sum lies in $e_1/2+L_2$, and
$Q(x)=t^2\lvert d\rvert/4\equiv-\tfrac14\bmod\Z$ there.
\end{proof}

\begin{lemma}[a level prime dividing both the conductor and $d_0$]\label{lem:ramlevel}
Let $p\,\|\,N$, $p\mid d_0$ and $p^k\,\|\,\ell$ with $k\ge1$. After the conjugations of
Lemma~\ref{lem:conductor} (and, for $p=2$, of Lemma~\ref{lem:conductor2}), $\lambda=\begin{pmatrix}0&b\\
p^2c'&0\end{pmatrix}$ with $b$ a unit and $bc'=d/p^2$ for odd $p$, and $\lambda=\begin{pmatrix}0&2b'\\4c'
&0\end{pmatrix}$ with $b'$ odd, $c'$ even and $b'c'=d/8$ for $p=2$. Then $L_{-,p}=\langle-1\rangle\perp
\langle-\lvert d\rvert/p^{\,2}\rangle$ for odd $p$ and $\langle-1\rangle\perp\langle-\lvert d\rvert/4\rangle$
for $p=2$: the anisotropic planes of Lemmas~\ref{lem:unimod}(4) and \ref{lem:unimod2}(1), with their
integral cosets $(0,r/p^k)$, $\rho=\ord_pr$, their series, and
\[
  \kappa^-_{\nu_r}(0)=-\frac{2(p^{\rho+1}-1)}{(p-1)p^k}\log p .
\]
The index $[L_p:L_{+,p}\oplus L_{-,p}]$ is $p^{2k}$, and the cosets of order $p$ ($\rho=k-1$), $p-1$
of them for odd $p$ and one for $p=2$, lie in $L_p^\vee$, over nonzero isotropic classes of
$\Ldual/L$ supported at $p$. So, unlike the fundamental ramified case (Remark~\ref{rem:ramified}),
the constant term enters: the $x=0$ pairs contribute $\tfrac{p^k-1}{p^k}\,\tfrac12c_\eta(0)\log p$ per
point, and the $x\ne0$ pairs the poles of $f$ at the CM points of $d/p^2,\dots$ as before.
\end{lemma}
\begin{proof}
The derivation of Lemma~\ref{lem:conductor}(1) used only $p^2\mid Q(\lambda)$ and optimality, so it
applies, with $a'^2+bc'=d/p^2$ now divisible by $p$; normalising $a'=0$ the plane is
$-\alpha^2-(p^2c'/b)\beta^2=-\alpha^2-(\lvert d\rvert/p^2)(1/b^2)\beta^2$ with $\ord_p(\lvert d\rvert/p^2)
=2k-1$ odd. For $p=2$ the derivation of Lemma~\ref{lem:conductor2}(1) gives $a'$ even (as $a'^2+2b'c'
=-\lvert d\rvert/4\equiv0\bmod4$), $b'$ odd and $c'$ even, and the conjugation by $\begin{pmatrix}1&0\\
2t&1\end{pmatrix}$ shifts $a'$ by $2tb'$, so $a'=0$; the plane is $-\alpha^2+(2c'/b')\beta^2$ with
$2c'/b'=-\lvert d\rvert/(4b'^2)=-4^k\lvert m\rvert \ell'^2/b'^2$, $d_0=4m$. The series and $\kappa^-$ are
those computed in Lemmas~\ref{lem:unimod}(4) and \ref{lem:unimod2}(1), which depend only on the plane.
The dual pairing is $\langle e_2,x\rangle=-p\gamma+(p^2c'/b)\beta$ (resp.\ $-2\gamma+(2c'/b')\beta$),
so $e_2/p^m\in L_p^\vee$ iff $m\le1$, and $Q(e_2/p)=-\lvert d\rvert/p^4\in\Z_p$ (resp.\ $Q(e_2/2)=
-\lvert d\rvert/16$) is integral since $k\ge1$; its class in $(\Ldual/L)_p$ is $e_{12}/p$ (resp.\ the
class of $e_{12}/2$), a nonzero isotropic vector. The index is the valuation of
$\det(\lambda,e_1,e_2)=2pc'$ (resp.\ of the determinant ratio $4^{2k}$ as in Lemma~\ref{lem:conductor2}).
The weight of the $x=0$ pairs is $-\tfrac14\cdot(p-1)\cdot\kappa^-_{\nu}(0)\cdot c_\eta(0)$ at $\rho=k-1$
for odd $p$, and $-\tfrac14\kappa^-_\nu(0)c_\eta(0)$ for $p=2$ where $p-1=1$.
\end{proof}

So at every prime of the conductor the term is determined, by Lemmas~\ref{lem:conductor},
\ref{lem:conductor2}, \ref{lem:unimod}, \ref{lem:unimod2} and \ref{lem:ramlevel}. \paragraph{Check.} One published value reaches this case: on $X_0^{21}(2)$, \cite{GY} give
$s(-16)=1$, with $-16=2^2\cdot(-4)$ and $2$ the level prime dividing $d_0$. There the term is $-4,-6,-8,-4
\log2$ on four of the covers and $0$ on the Hauptmodul forms (for them the $x=0$ pair, with
$\tfrac12c_\eta(0)$, is cancelled by the pole of $f$ at $\tau_{-4}$ through $x=\pm\lambda_0/2$), so
the Hauptmodul row cannot see it; but every form is a function of $s$ with known divisor, supported
at $\tau_{-4}$ (the pole, where $s=\infty$), $\tau_{-7}$, $\tau_{-28}$, $\tau_{-84}$ and $\tau_{-168}$,
all with published $s$. Fitting each form's constant at one fundamental point and checking it at ten
to thirteen others, the value at $-16$ forced by the divisor agrees with the pipeline's for all nine
forms; without the term the four covers are off by $2^4,2^6,2^8,2^4$ (\texttt{level-p2/ramcheck.py}). \paragraph{Check.} Here the published tables do reach the case, once the points are forced in:
on the odd-level bases $X_0^{35}(1)$, $X_0^{39}(1)$, $X_0^{51}(1)$, $X_0^{55}(1)$ and $X_0^{57}(1)$,
\cite{GY} list $s$ at $d=-28$ and $-60$ ($2$ split in $k$), $-12$ ($2$ inert), and $-16=2^2\cdot(-4)$
($d_0$ even), and at $-315=3^2\cdot(-35)$ on $X_0^{35}(1)$ ($3\mid d_0$). The terms are
$\log2$ on the Hauptmodul form at the split points, $-\log2$ at $-12$ and $-16$ (and $-2,-4,-6,-8\log2$
on the covers), $\log3$ at $-315$, and with them the pipeline reproduces every one of these eight
published values, while every other row of those tables is unchanged
(\texttt{tests/\_offline/GuoYang\_*.m}, which now force these points). Without a term at $2$ the
Hauptmodul value is off by a factor $2$ at each of them.

\section{Computing the constant terms exactly}
\begin{remark}[computing $c_\eta(0)$ exactly]\label{rem:exact}
The constant term $c_\eta(0)=\sum_w\rho(w^{-1})e_0[\eta]\,a_0(f|w)$, a sum over the cosets $w$ of
$\Gamma_0(M)\backslash SL_2(\mathbf Z)$, is an algebraic number, and it can be computed with no
transcendental step. The per-coset product is well defined on $\Gamma_0(M)\backslash SL_2(\mathbf Z)$
(both factors change by the character of $f$ on $\Gamma_0(M)$ under left multiplication, and $e_0$ is
an eigenvector for that character), and it is constant on each class $g=\gcd(c,M)$ of the bottom row
of $w$ for two reasons. Right multiplication by $T$ leaves $a_0(f|w)$ unchanged and multiplies the
$\eta$-component of $\rho(w^{-1})e_0$ by $\mathbf e(\pm Q(\eta))=1$, $Q(\eta)$ being an integer at
the isotropic $\eta$ in question, so the product is a function of the cusp $\Gamma_0(M)w\infty$ alone;
and the cusps of $\Gamma_0(M)$ with $\gcd(c,M)=g$ number $\varphi(\gcd(g,M/g))$, which is $1$ for
every $g$ because $M$ is four times an odd squarefree number, so each class is a single cusp. One
representative $\gamma=\bigl(\begin{smallmatrix}a&b\\c&d\end{smallmatrix}\bigr)$ with $c>0$ per
class therefore suffices, both factors taken for the lift $(\gamma,\sqrt{c\tau+d})$ of the metaplectic
group; the check at a second representative, which this makes redundant, is off by default in the
pipeline. The first factor is Shintani's formula for the Weil representation at a general element
\cite[Prop.~1.6]{Shintani}, in the form given by Scheithauer \cite[Thm.~4.7]{Scheithauer} and
Str\"omberg \cite{Stromberg}: for $\gamma$ with $c>0$, $\rho_L(\gamma)e_\delta$ is a Gauss sum over
$L/cL$ with the prefactor $\mathbf e\bigl((b^--b^+)/8\bigr)\,c^{-n/2}\,\lvert\Ldual/L\rvert^{-1/2}$,
$(b^+,b^-)$ the signature and $n=b^++b^-$, the same formula one gets by writing
$\gamma\tau=a/c-1/(c(c\tau+d))$ and applying Poisson summation to the inversion in the theta
transformation formula \cite[Thm.~4.1]{Borcherds}. Two conventions enter. The input forms transform
under the dual of Borcherds' $\rho_L$, its complex conjugate: the coset sum $\sum_w(f|w)\rho(w^{-1})e_0$
is well defined only if $\rho(Z^{-1})e_0=\mathbf e(\tfrac14)e_0$ for $Z=S^2$, and at signature $(1,2)$
that is the dual's sign ($\rho_L(Z)e_\gamma=\mathbf e(\tfrac14)e_{-\gamma}$ for Borcherds' $\rho_L$).
And what is needed is $\rho(\gamma^{-1})$, whose lower-left entry is $-c<0$; writing
$\gamma^{-1}=Z^{-1}\cdot(-\gamma^{-1})$ and applying the formula to $-\gamma^{-1}$ gives the dual's
prefactor $\mathbf e(-\tfrac18)$ times $\mathbf e(\tfrac14)$ from $Z^{-1}$. So, for the dual
representation $\rho$ of a lattice of signature $(1,2)$,
\[
\rho(\gamma^{-1})e_0[\eta]=\mathbf e(\tfrac18)\,c^{-3/2}\,\lvert\Ldual/L\rvert^{-1/2}\,
\mathbf e\Bigl(\frac{d\,Q(\eta)}c\Bigr)\sum_{\nu\in L/cL}\mathbf e\Bigl(\frac{a\,Q(\nu)+(\eta,\nu)}c\Bigr),
\]
the Gauss sum factoring over the prime powers of $c$ (the quadratic term takes the cofactor, the
linear term does not); the entry $d$ in the prefactor and $a$ in the sum are not interchangeable.
The second factor is the Dedekind-eta multiplier system: for a monomial $\prod_d\eta(d\tau)^{r_d}$
of weight $\tfrac12$, write $\bigl(\begin{smallmatrix}d&0\\0&1\end{smallmatrix}\bigr)\gamma=g_d
\bigl(\begin{smallmatrix}a_d&b_d\\0&e_d\end{smallmatrix}\bigr)$ with $g_d\in SL_2(\mathbf Z)$ of
positive lower-left entry; then Apostol's transformation law $\eta(g\tau)=\varepsilon(g)\,
(-i(c\tau+d))^{1/2}\eta(\tau)$, with $\varepsilon(g)=\exp\bigl(\pi i\bigl(\frac{a+d}{12c}+s(-d,c)
\bigr)\bigr)$ a $24$th root of unity given by a Dedekind sum, gives
\[
\Bigl(\prod_d\eta(d\tau)^{r_d}\Bigr)\Big|\gamma
=\mathbf e(-\tfrac18)\prod_d\bigl[\varepsilon(g_d)\,\mathbf e\bigl(\tfrac{b_d}{24e_d}\bigr)\,
e_d^{-1/2}\bigr]^{r_d}\;q^{\sum_d r_da_d/(24e_d)}\prod_d\prod_{n\ge1}
\bigl(1-\mathbf e(nb_d/e_d)\,q^{na_d/e_d}\bigr)^{r_d},
\]
whose constant term is read off: the Borcherds inputs are holomorphic at every cusp other than
$\infty$ and $0$ and so are their monomials (lattice points of the polytope), so at the classes
$1<g<M$ the constant term is the product of the roots of unity when the order $\sum_dr_da_d/e_d$
vanishes and $0$ otherwise, and at the cusp $0$ the representative $S$ makes the $q$-series rational.
Every root of unity lies in $\mu_{24M}$ and the square roots are those of integers dividing $2M$, so
the whole computation takes place in $\mathbf Q(\zeta_{24M})$ --- in practice in the group ring
$\mathbf Q[x]/(x^{24M}-1)$ with the square roots as Gauss sums, reduced modulo the cyclotomic
polynomial at the end, where the result must be a rational number. This is what the routine
\texttt{M0MultipliersAlgebraic} does, and what the pipeline now uses. It reproduces the measured
ground truth on $X_0^{15}(2)$ and the numerical coset evaluation on $X_0^{21}(2)$, $X_0^{10}(3)$,
$X_0^{22}(3)$, $X_0^{34}(11)$, on the four probe eta quotients of Remark~\ref{rem:ceta}
(all three classes), and on the Borcherds inputs of the three composite-level bases $X_0^6(35)$,
$X_0^{10}(21)$, $X_0^{14}(15)$ that have since been built, in seconds to minutes where the numerical
route took an hour or more.
\end{remark}

\begin{thebibliography}{9}
\bibitem{Schofer} J.~Schofer, \emph{Borcherds forms and generalizations of singular moduli},
J.\ reine angew.\ Math.\ \textbf{629} (2009), 1--36.
\bibitem{KY} S.~Kudla and T.~Yang, \emph{Eisenstein series for $SL(2)$},
Sci.\ China Math.\ \textbf{53} (2010), 2275--2316.
\bibitem{Yang} Y.~Yang, \emph{Special values of hypergeometric functions and periods of CM elliptic
curves}, arXiv:1503.07971.
\bibitem{GY} J.-W.~Guo and Y.~Yang, \emph{Equations of hyperelliptic Shimura curves},
Compos.\ Math.\ \textbf{153} (2017), 1--40.
\bibitem{Err} E.~Errthum, \emph{Singular moduli of Shimura curves}, Canad.\ J.\ Math.\ \textbf{63}
(2011), 826--861.
\bibitem{BY} J.~H.~Bruinier and T.~Yang, \emph{Faltings heights of CM cycles and derivatives of
$L$-functions}, Invent.\ Math.\ \textbf{177} (2009), 631--681.
\bibitem{BEY} J.~H.~Bruinier, S.~Ehlen and T.~Yang, \emph{CM values of higher automorphic Green
functions for orthogonal groups}, Invent.\ Math.\ \textbf{225} (2021), 693--785.
\bibitem{Borcherds} R.~E.~Borcherds, \emph{Automorphic forms with singularities on Grassmannians},
Invent.\ Math.\ \textbf{132} (1998), 491--562.
\bibitem{Shintani} T.~Shintani, \emph{On construction of holomorphic cusp forms of half integral
weight}, Nagoya Math.\ J.\ \textbf{58} (1975), 83--126.
\bibitem{Scheithauer} N.~R.~Scheithauer, \emph{The Weil representation of $SL_2(\mathbf Z)$ and some
applications}, Int.\ Math.\ Res.\ Not.\ IMRN (2009), no.~8, 1488--1545.
\bibitem{Stromberg} F.~Str\"omberg, \emph{Weil representations associated with finite quadratic
modules}, Math.\ Z.\ \textbf{275} (2013), 509--527.
\end{thebibliography}

\end{document}
